% ====== 第五部分 专利申请及事务处理 ======
\cleardoublepage\setguidefootmark{5}{第五部分}\resetchapters
\thispagestyle{plainempty}\hypertarget{gpos.p5}{}\begingroup\centering\guidecenterpage
\dadati\fontsize{29.04bp}{43bp}\selectfont
\vspace*{242bp}
第五部分\par
专利申请及事务处理\par
\phantomsection
\addcontentsline{toc}{part}{第五部分\quad 专利申请及事务处理}
\endgroup\clearpage
\cleardoublepage\pagestyle{tocpage}
\guideoutline[1]{分目录}{parttoc5}
\gdef\guidetabcontent{\hbox{第}\hbox{五}\hbox{部}\hbox{分}}
{\centering\guidecenterpage\heiti\fontsize{14.04bp}{21bp}\selectfont 目\quad 录\par}
\vspace{6bp}
\printparttoc{p5}
\clearpage\pagestyle{main}
\startparttoc{p5}
\guidechapter{gpos.5.1}{第一章}{专利申请文件及手续}{}{0bp}
\guidesection{gpos.5.1.1}{1.}{引言}{}{74.4bp}
申请人就一项发明创造要求获得专利权的，应当根据专利法及其实施细则的规定向专利局提出专利申请。在专利审批程序中，申请人根据专利法及其实施细则的规定或者审查员的要求，还需要办理各种与该专利申请有关的事务。申请人向专利局提出专利申请以及在专利审批程序中办理其他专利事务，统称为专利申请手续。

申请人提出专利申请，向专利局提交的专利法第二十六条规定的请求书、说明书、权利要求书、说明书附图和摘要或者专利法第二十七条规定的请求书、图片或者照片、简要说明等文件，称为专利申请文件；在提出专利申请的同时或者提出专利申请之后，申请人（或者专利权人）、其他相关当事人在办理与该专利申请（或者专利）有关的各种手续时，提交的除专利申请文件以外的各种请求、申报、意见陈述、补正以及各种证明、证据材料，称为其他文件。

办理各种手续应当提交相应的文件，缴纳相应的费用，并且符合相应的期限要求。

\guidesection{gpos.5.1.2}{2.}{办理专利申请手续的形式}{\guidenote{细则 2}}{74.4bp}
专利申请手续应当以符合规定的电子形式、纸件形式等书面形式办理。

以口头、电话、实物、传真、电子邮件等形式办理的，视为未提出，不产生法律效力，另有规定的除外。

\guidesubsection{gpos.5.1.2.1}{2.1}{电子形式}{}{74.4bp}
申请人以电子形式提出专利申请并被受理的，在审批程序中应当通过专利局指定的特定电子系统以电子形式提交相关文件，另有规定的除外。不符合规定的，该文件视为未提交。

\guidesubsection{gpos.5.1.2.2}{2.2}{纸件形式}{}{74.4bp}
申请人以纸件形式提出专利申请并被受理的，在审批程序中应当以纸件形式提交相关文件。除另有规定外，申请人以电子形式提交的相关文件视为未提交。

对受理的纸件专利申请文件和其他文件，专利局进行扫描并存入数据库。以纸件形式提交的专利申请文件及其他文件，经专利局转换为电子形式文件并记录在电子系统数据库中，具有与原纸件形式文件同等的效力。

\guidesubsection{gpos.5.1.2.3}{2.3}{纸件申请和电子申请的转换}{}{74.4bp}
申请人、复审请求人或者专利代理机构可以请求将纸件申请转换为电子申请，涉及国家安全或者重大利益需要保密的专利申请除外。

提出请求的申请人、复审请求人或者专利代理机构应当通过专利局指定的特定电子系统提出请求。使用其他方式提出请求的，该请求视为未提出。

\guidesection{gpos.5.1.3}{3.}{适用文字}{}{102.0bp}
\guidesubsection{gpos.5.1.3.1}{3.1}{中文}{\guidenote{细则 3}}{0bp}
专利申请文件以及其他文件，除由外国政府部门出具的或者在外国形成的证明或者证据材料外，应当使用中文。

审查员以申请人提交的中文专利申请文本为审查的依据。申请人在提出专利申请的同时提交的外文申请文本，供审查员在审查程序中参考，不具有法律效力。

\guidesubsection{gpos.5.1.3.2}{3.2}{汉字}{}{74.4bp}
本章第3.1节中的“中文”一词是指汉字。专利申请文件及其他文件应当使用汉字，词、句应当符合现代汉语规范。

汉字应当以国家公布的简化字为准。申请文件中的异体字、繁体字、非规范简化字，审查员可以依职权予以改正或者通知申请人补正。

\guidesubsection{gpos.5.1.3.3}{3.3}{外文的翻译}{}{74.4bp}
\guidenote{细则 3.1}专利申请文件是外文的，应当翻译成中文，其中外文科技术语应当按照规定译成中文，并采用规范用语。外文科技术语没有统一中文译法的，可按照一般惯例译成中文，并在译文后的括号内注明原文。计量单位应当使用国家法定计量单位，包括国际单位制计量单位和国家选定的其他计量单位，必要时可以在括号内同时标注本领域公知的其他计量单位。

\guidenote{细则 3.2}当事人在提交外文证明文件、证据材料（例如优先权证明文本、转让证明等）时，应当同时附具中文题录译文，审查员认为必要时，可以要求当事人在规定的期限内提交全文中文译文或者摘要中文译文；期满未提交译文的，视为未提交该文件。

\guidesection{gpos.5.1.4}{4.}{标准表格}{}{74.4bp}
办理专利申请（或者专利）手续时应当使用专利局制定的标准表格。标准表格由专利局按照一定的格式和样式统一制定、修订和公布。

办理专利申请（或者专利）手续时以非标准表格提交的文件，审查员可以根据有关规定发出补正通知书或者针对该手续发出视为未提出通知书。

但是，申请人在答复补正通知书或者审查意见通知书时，提交的补正书或者意见陈述书为非标准格式的，只要写明申请号，表明是对申请文件的补正，并且签字或者盖章符合规定的，可视为文件格式符合要求。

\guidesubsection{gpos.5.1.4.1}{4.1}{纸张}{}{74.4bp}
各种文件使用的纸张应当柔韧、结实、耐久、光滑、无光、白色。其质量应当与80克胶版纸相当或者更高。

\guidesubsection{gpos.5.1.4.2}{4.2}{规格}{}{74.4bp}
说明书、说明书附图、权利要求书、说明书摘要、摘要附图、图片或者照片、简要说明与其他表格用纸的规格均应为297毫米×210毫米（A4）。

\guidesubsection{gpos.5.1.4.3}{4.3}{页边}{}{74.4bp}
申请文件的顶部（有标题的，从标题上沿至页边）应当留有25毫米空白，左侧应当留有25毫米空白，右侧应当留有15毫米空白，底部从页码下沿至页边应当留有15毫米空白。

\guidesection{gpos.5.1.5}{5.}{书写规则}{}{102.0bp}
\guidesubsection{gpos.5.1.5.1}{5.1}{打字或者印刷}{}{0bp}
请求书、权利要求书、说明书、说明书摘要、说明书附图和摘要附图中文字部分以及简要说明应当打字或者印刷。上述文件中的数学式和化学式可以按照制图方式手工书写。

其他文件除另有规定外，可以手工书写，但字体应当工整，不得涂改。

\guidesubsection{gpos.5.1.5.2}{5.2}{字体及规格}{}{74.4bp}
各种文件应当使用宋体、仿宋体或者楷体，不得使用草体或者其它字体。

纸件申请的字高应当不低于3.5毫米，行距应当在2.5毫米至3.5毫米之间。

\guidesubsection{gpos.5.1.5.3}{5.3}{书写方式}{}{74.4bp}
各种文件除另有规定外，应当单面、纵向使用。自左至右横向书写，不得分栏书写。

一份文件不得涉及两件以上专利申请（或者专利），一页纸上不得包含两种以上文件（例如一页纸不得同时包含说明书和权利要求书）。

\guidesubsection{gpos.5.1.5.4}{5.4}{书写内容}{}{74.4bp}
文件各栏目应当如实、详尽填写，同一内容在不同栏目或者不同文件中应当填写一致。例如地址栏目应当按照行政区划填写完整，邮政编码与地址相符；申请人的签字或者盖章应当与申请人栏目中填写的内容一致。

\guidesubsection{gpos.5.1.5.5}{5.5}{字体颜色}{}{74.4bp}
字体颜色应当为黑色，字迹应当清晰、牢固、不易擦、不褪色，以能够满足复印、扫描的要求为准。

\guidesubsection{gpos.5.1.5.6}{5.6}{编写页码}{}{74.4bp}
各种文件应当分别用阿拉伯数字顺序编写页码。页码应当置于每页下部页边的上沿，并左右居中。

\guidesection{gpos.5.1.6}{6.}{证明文件}{}{74.4bp}
专利申请审批程序中常用的证明文件有国籍证明、注册地证明、优先权证明（在先申请文件副本）、优先权转让证明、生物材料样品保藏证明、申请人（或者专利权人）名称变更或者权利转移证明、文件寄发日期证明等。

各种证明文件应当由有关主管部门出具或者由当事人签署。各种证明文件应当提供原件；证明文件是复印件的，应当经公证或者由主管部门加盖公章予以确认。申请人办理专利电子申请的相关手续时，对专利法及其实施细则和本指南中规定的应当以原件形式提交的相关证明文件，可以提交原件的电子扫描文件。专利局认为必要时，可以要求申请人在指定期限内提交原件。

申请人可以将证明文件原件在专利局存档备案，在办理相关手续时应当注明证明文件备案号。

申请人（或者专利权人）、其他相关当事人在办理与专利申请（或者专利）有关的各种手续时，应当遵循诚实信用原则。对于涉及违反诚实信用原则的专利申请（或者专利），专利局不予批准相关手续，已经批准的，予以撤销。

\guidesection{gpos.5.1.7}{7.}{文件份数}{}{74.4bp}
申请人向专利局提交的申请文件及其他文件，例如专利代理委托书、实质审查请求书、著录项目变更申报书、转让合同等，应当为一份，另有规定的除外。

\guidesection{gpos.5.1.8}{8.}{签字或者盖章}{}{74.4bp}
\guidenote{细则 146.1}向专利局提交的专利申请文件或者其他文件，应当按照规定签字或者盖章。其中未委托专利代理机构的申请，应当由申请人（或者专利权人）、其他利害关系人或者其代表人签字或者盖章，办理直接涉及共有权利的手续，应当由全体权利人签字或者盖章；委托了专利代理机构的，应当由专利代理机构盖章，必要时还应当由申请人（或者专利权人）、其他利害关系人或者其代表人签字或者盖章。

\guidesection{gpos.5.1.9}{9.}{专利业务办理系统的注册}{}{74.4bp}
专利业务办理系统是专利局指定的特定电子系统，是可以集中实现专利电子申请、电子缴费、集成电路布图设计电子申请等多个业务的网上业务办理平台。

专利业务办理系统注册用户是指已经与国家知识产权局签订了专利业务办理系统用户服务协议，实名办理了用户注册手续的自然人、法人和非法人组织。

注册用户的注册方式、注册信息维护、数字证书的使用等手续，应当按照用户服务协议的内容办理。

申请人有两人以上且未委托专利代理机构，并以电子形式提交专利申请文件及其他文件的，应当由代表人提交。

\guidechapter{gpos.5.2}{第二章}{专利费用}{}{0bp}
\guidesection{gpos.5.2.1}{1.}{费用缴纳的期限}{}{74.4bp}
{\setlength{\parindent}{17.1bp}\guidenote{细则 112}（1）申请费的缴纳期限是自申请日起两个月内，或者自收到受理通知书之日起15日内。需要在该期限内缴纳的费用有优先权要求费和申请附加费以及发明专利申请的公布印刷费。\par}

优先权要求费是指申请人要求外国优先权或者本国优先权时，需要缴纳的费用，该项费用的数额以作为优先权基础的在先申请的项数计算。

申请附加费是指申请文件的说明书（包括附图、序列表）页数超过30页或者权利要求超过10项时需要缴纳的费用，该项费用的数额以页数或者项数计算。

公布印刷费是指发明专利申请公布需要缴纳的费用。

未在规定的期限内缴纳或者缴足申请费（含公布印刷费、申请附加费）的，该申请被视为撤回。未在规定的期限内缴纳或者缴足优先权要求费的，视为未要求优先权。

{\setlength{\parindent}{17.1bp}\guidenote{法 35.1}\guidenote[15.6bp]{细则 113}（2）实质审查费的缴纳期限是自申请日（有优先权要求的，自最早的优先权日）起三年内。该项费用仅适用于发明专利申请。\par}

{\setlength{\parindent}{17.1bp}\guidenote{细则 116.2}（3）延长期限请求费的缴纳期限是在相应期限届满之日前。该项费用以要求延长的期限长短（以月为单位）计算。\par}

{\setlength{\parindent}{17.1bp}\guidenote{细则 6 及 116.1}（4）恢复权利请求费的缴纳期限是自当事人收到专利局确认权利丧失通知之日起两个月内。\par}

{\setlength{\parindent}{17.1bp}\guidenote{法 41.1}\guidenote[15.7bp]{细则 113}（5）复审费的缴纳期限是自申请人收到专利局作出的驳回决定之日起三个月内。\par}

{\setlength{\parindent}{17.1bp}\guidenote{细则 60 及 114}（6）授权当年年费的缴纳期限是自申请人收到专利局作出的授予专利权通知书和办理登记手续通知书之日起两个月内。\par}

{\setlength{\parindent}{17.1bp}\guidenote{细则 115}（7）年费及其滞纳金的缴纳期限参照本部分第九章第4.2.1节的规定。\par}

{\setlength{\parindent}{17.1bp}\guidenote{细则 116.3}（8）著录事项变更费、专利权评价报告请求费、无效宣告请求费的缴纳期限是自提出相应请求之日起一个月内。\par}

\guidesection{gpos.5.2.2}{2.}{费用支付和结算方式}{\guidenote{细则 111.1}}{74.4bp}
费用可以直接向专利局（包括专利局各代办处）缴纳，也可以通过邮局或者银行汇付，或者以规定的其他方式缴纳。专利局代办处的收费范围另行规定。

费用通过邮局或者银行汇付的，应当在汇单上写明正确的申请号（或者专利号）以及缴纳的费用名称，且不得设置取款密码。不符合上述规定的，视为未办理缴费手续。

在汇单上还应当写明汇款人姓名或者名称及其通讯地址（包括邮政编码）。同一专利申请（或者专利）缴纳的费用为两项以上的，应当分别注明每项费用的名称和金额，并且各项费用的金额之和应当等于缴纳费用的总额。

同一汇单中包括多个专利申请（或者专利），其缴纳费用的总额少于各项专利申请（或者专利）费用金额之和的，处理方法如下：

{\setlength{\parindent}{17.1bp}（1）缴费人对申请号（或者专利号）标注顺序号的，按照标注的顺序分割费用；\par}

{\setlength{\parindent}{17.1bp}（2）缴费人未对申请号（或者专利号）标注顺序号的，按照从左至右、从上至下的顺序分割费用。\par}

造成其中部分专利申请（或者专利）费用金额不足或者无费用的，视为未办理缴费手续。

\guidenote{细则 111.2}费用通过邮局汇付，且在汇单上写明申请号（或者专利号）以及费用名称的，以邮局取款通知单上的汇出日为缴费日。邮局取款通知单上的汇出日与中国邮政普通汇款收据上收汇邮戳日表明的日期不一致的，以当事人提交的中国邮政普通汇款收据原件或者经公证的收据复印件上表明的收汇邮戳日为缴费日。审查员认为当事人提交的证据有疑义时，可以要求当事人提交汇款邮局出具的加盖部门公章的证明材料。

费用通过银行汇付，且写明申请号（或者专利号）以及费用名称的，以银行实际汇出日为缴费日。当事人对缴费日有异议，并提交银行出具的加盖部门公章的证明材料的，以证明材料确认的汇出日重新确定缴费日。

费用通过邮局或者银行汇付，未写明申请号（或者专利号）的，费用退回。费用退回的，视为未办理缴费手续。

因缴费人信息填写不完整或者不准确，造成费用不能退回或者退款无人接收的，费用暂时存入专利局账户（以下简称暂存）。费用入暂存的，视为未办理缴费手续。

各种费用以人民币结算。使用外币支付的费用，按照银行将该费用结汇至专利局账户之日的汇率结算，以结算之日为缴费日。

\guidesection{gpos.5.2.3}{3.}{收费的减缴}{}{74.4bp}
\guidenote{细则 117}申请人（或者专利权人）缴纳专利费用有困难的，或者其专利处于开放许可实施期间的，可以根据专利收费减缴办法向专利局提出收费减缴的请求。

\guidesubsection{gpos.5.2.3.1}{3.1}{可以减缴的费用种类}{}{74.4bp}
{\setlength{\parindent}{17.1bp}（1）申请费（不包括公布印刷费、申请附加费）；\par}

{\setlength{\parindent}{17.1bp}（2）发明专利申请实质审查费；\par}

{\setlength{\parindent}{17.1bp}（3）复审费；\par}

{\setlength{\parindent}{17.1bp}（4）年费（自授予专利权当年起十年的年费、开放许可实施期间的年费）。\par}

\guidesubsection{gpos.5.2.3.2}{3.2}{收费减缴的手续}{}{74.4bp}
提出专利申请时以及在审批程序中，申请人（或者专利权人）只能请求减缴应当缴纳但尚未到期的费用。

请求减缴专利收费的，应当提出收费减缴请求，并在提出请求前提前办理专利费减备案手续。办理专利开放许可实施合同备案的，视为提出年费减缴请求，无需办理专利费减备案手续。

专利收费减缴请求手续应当由申请人（或者专利权人）或者其代表人办理；已委托专利代理机构的，应当由专利代理机构办理。

收费减缴请求符合规定的，审查员应当发出收费减缴审批通知书，同时注明收费减缴的比例和种类。收费减缴请求不符合规定的，审查员应当发出收费减缴审批通知书，并说明不予减缴的理由。

专利收费减缴办法另行规定。

\guidesection{gpos.5.2.4}{4.}{费用的暂存与退款}{}{102.0bp}
\guidesubsection{gpos.5.2.4.1}{4.1}{暂存}{}{0bp}
由于费用汇单字迹不清或者缺少必要事项造成既不能开出收据又不能退款的，应当将该款项暂存在专利局账户上。经缴款人提供证明后，对于能够查清内容的，应当及时开出收据或者予以退款。开出收据的，以出暂存之日为缴费日。但是，对于自收到专利局关于权利丧失的通知之日起两个月内向专利局提交了证据，表明是由于银行或者邮局原因导致汇款暂存的，应当以原汇出日为缴费日。暂存满三年仍无法查清其内容的，进行清账上缴。

\guidesubsection{gpos.5.2.4.2}{4.2}{退款}{}{102.0bp}
\guidesubsubsection{gpos.5.2.4.2.1}{4.2.1}{退款的原则}{}{0bp}
\guidenote{细则 111.3}多缴、重缴、错缴专利费用的，当事人可以自缴费日起三年内，提出退款请求。

符合国务院发展改革部门、财政部门及国务院专利行政部门发布的公告和通知的有关规定的，当事人可以提出退款请求。

符合规定的，专利局应当予以退款。

\guideparagraph{gpos.5.2.4.2.1.1}{4.2.1.1}{当事人可以请求退款的情形}{}{74.4bp}
{\setlength{\parindent}{17.1bp}（1）多缴费用的情形：如当事人应当缴纳年费为600元，在规定的期限内实际缴纳费用为650元，可以对多缴的50元提出退款请求。\par}

{\setlength{\parindent}{17.1bp}（2）重缴费用的情形：如提出一次著录项目变更请求应当缴纳著录事项变更费200元，当事人缴纳200元后，再次缴纳了200元，当事人可以对再次缴纳的200元提出退款请求。\par}

{\setlength{\parindent}{17.1bp}（3）错缴费用的情形：如当事人缴费时写错费用种类、申请号（或者专利号）的；或者因缴费不足、逾期缴费导致权利丧失的，或者权利丧失后缴纳专利费用的，当事人可以提出退款请求。\par}

{\setlength{\parindent}{17.1bp}（4）当事人对于进入实质审查阶段的发明专利申请，在第一次审查意见通知书答复期限届满前主动申请撤回的，可以请求退还50\%的发明专利申请实质审查费，已提交答复意见的除外。\par}

\guideparagraph{gpos.5.2.4.2.1.2}{4.2.1.2}{专利局主动退款的情形}{}{74.4bp}
下列情形一经核实，专利局应当主动退款。

{\setlength{\parindent}{17.1bp}（1）专利申请已被视为撤回或者撤回专利申请的声明已被批准后，并且在专利局作出发明专利申请进入实质审查阶段通知书之前，已缴纳的实质审查费。\par}

{\setlength{\parindent}{17.1bp}（2）在专利权终止或者宣告专利权全部无效的决定公告后缴纳的年费。\par}

{\setlength{\parindent}{17.1bp}（3）恢复权利请求审批程序启动后，专利局作出不予恢复权利决定的，当事人已缴纳的恢复权利请求费及相关费用。\par}

\guideparagraph{gpos.5.2.4.2.1.3}{4.2.1.3}{不予退款的情形}{}{74.4bp}
{\setlength{\parindent}{17.1bp}（1）对多缴、重缴、错缴的费用，当事人在自缴费日起三年后才提出退款请求的。\par}

{\setlength{\parindent}{17.1bp}（2）当事人不能提供错缴费用证据的。\par}

{\setlength{\parindent}{17.1bp}（3）在收费减缴请求被批准之前已经按照规定缴纳的各种费用，当事人又请求退款的。\par}

\guidesubsubsection{gpos.5.2.4.2.2}{4.2.2}{退款的手续}{}{102.0bp}
\guideparagraph{gpos.5.2.4.2.2.1}{4.2.2.1}{退款请求的提出}{}{0bp}
退款请求人应当是该款项的缴款人。申请人（或者专利权人）、专利代理机构作为非缴款人请求退款的，应当声明是受缴款人委托办理退款手续。

请求退款应当书面提出、说明理由并附具相应证明，例如，专利局开出的费用收据复印件、邮局或者银行出具的汇款凭证等。提供邮局或者银行的证明应当是原件，不能提供原件的，应当提供经出具部门加盖公章确认的或者经公证的复印件。

退款请求应当注明申请号（或者专利号）和要求退款的款项的信息（如票据号、费用金额等）及收款人信息。当事人要求通过邮局退款的，收款人信息包括姓名、地址和邮政编码；当事人要求通过银行退款的，收款人信息包括姓名或者名称、开户行、账号等信息。

\guideparagraph{gpos.5.2.4.2.2.2}{4.2.2.2}{退款的处理}{}{74.4bp}
经核实可以退款的，专利局应当按照退款请求中注明的收款人信息退款。

退款请求中未注明收款人信息的，退款请求人是申请人（或者专利权人）或者专利代理机构的，应当按照文档中记载的相应的地址和姓名或者名称退款。

完成退款处理后，审查员应当发出退款审批通知书。经核实不予退款的，审查员应当在退款审批通知书中说明不予退款的理由。

\guidesubsubsection{gpos.5.2.4.2.3}{4.2.3}{退款的效力}{}{74.4bp}
被退的款项视为自始未缴纳。

\guidesubsubsection{gpos.5.2.4.2.4}{4.2.4}{特殊情形的处理}{}{102.0bp}
\guideparagraph{gpos.5.2.4.2.4.1}{4.2.4.1}{因银行或者邮局责任造成必要缴费信息不全被退款的情形}{}{0bp}
因银行或者邮局责任造成必要缴费信息（如申请号、费用名称等）不完整被退款，当事人提出异议的，应当以书面形式陈述意见，并附具汇款银行或者邮局出具的加盖公章的证明。该证明至少应当包括：汇款人姓名或者名称、汇款金额、汇款日期、汇款时所提供的申请号（或者专利号）、费用名称等内容。同时当事人应当重新缴纳已被退回的款项。

符合上述规定的，原缴费日视为重新缴纳款项的缴费日，因此导致已作出的处分决定需要更改的，审查员应当发出修改更正通知书。不符合上述规定的，审查员应当发出通知书通知当事人，该款项视为未缴纳。

\guideparagraph{gpos.5.2.4.2.4.2}{4.2.4.2}{因汇款人汇款后又取回汇款造成汇单无法兑付的情形}{}{74.4bp}
经邮局确认汇款人的汇款已被取回无法兑付的，该款项视为未缴纳。

\guideparagraph{gpos.5.2.4.2.4.3}{4.2.4.3}{通知书中涉及专利费用的信息存在错误的情形}{}{74.4bp}
通知书中涉及专利费用信息存在错误的，专利局应予更正。当事人应当按照专利局更正后的费用信息办理费用相关事宜。

\guidesection{gpos.5.2.5}{5.}{费用的查询}{}{74.4bp}
当事人需要查询费用缴纳情况的，应当提供银行汇单复印件或者邮局汇款凭证复印件（未收到专利局收费收据的）或者提供收据复印件（已收到专利局收费收据的）。查询时效为一年，自汇出费用之日起算。

\guidesection{gpos.5.2.6}{6.}{费用种类的转换}{}{74.4bp}
对于同一专利申请（或专利）缴纳费用时，费用种类填写错误的，缴纳该款项的当事人可以在转换后费用的缴纳期限内提出转换费用种类请求并附具相应证明，经专利局确认后可以对费用种类进行转换。但不同申请号（或者专利号）之间的费用不能转换。

当事人缴纳的费用种类明显错误，审查员可以依职权对费用种类进行转换。依职权转换费用种类的，应当通知当事人。

费用种类转换的，缴费日不变。

\guidesection{gpos.5.2.7}{7.}{缴费信息的补充}{}{74.4bp}
费用通过邮局或者银行汇付时遗漏必要缴费信息的，应当在汇款当日通过专利局规定的方式及要求补充。当日补充不完整而再次补充的，以专利局收到完整缴费信息之日为缴费日。

\guidechapter{gpos.5.3}{第三章}{受理}{}{0bp}
\guidesection{gpos.5.3.1}{1.}{受理部门}{}{74.4bp}
专利局受理部门包括专利局初审及流程管理部受理处、复审和无效审理部立案及流程管理处和专利局各代办处。初审及流程管理部受理处负责受理专利申请及其他有关文件，复审和无效审理部立案及流程管理处受理与复审和无效宣告请求有关的文件，各代办处按照相关规定受理专利申请及其他有关文件。

专利局受理部门应当开设受理窗口。未经过受理登记的文件，不得进入审批程序。

专利局受理部门的地址由专利局以公告形式公布。邮寄或者直接交给专利局的任何个人或者非受理部门的申请文件和其他有关文件，其邮寄文件的邮戳日或者提交文件的提交日都不具有确定申请日和递交日的效力。

\guidesection{gpos.5.3.2}{2.}{专利申请的受理与不受理}{}{102.0bp}
\guidesubsection{gpos.5.3.2.1}{2.1}{受理条件}{}{0bp}
\guidenote{细则 44}专利申请符合下列条件的，专利局应当受理：

{\setlength{\parindent}{17.1bp}（1）申请文件中有请求书。该请求书中申请专利的类别明确；写明了申请人姓名或者名称及其地址。\par}

{\setlength{\parindent}{17.1bp}（2）发明专利申请文件中有说明书和权利要求书；实用新型专利申请文件中有说明书、说明书附图和权利要求书；外观设计专利申请文件中有图片或者照片和简要说明。\par}

{\setlength{\parindent}{17.1bp}（3）申请文件是使用中文打字或者印刷的。全部申请文件的字迹和线条清晰可辨，没有涂改，能够分辨其内容。发明或者实用新型专利申请的说明书附图和外观设计专利申请的图片是用不易擦去的笔迹绘制，并且没有涂改。\par}

{\setlength{\parindent}{17.1bp}（4）申请人是外国人、外国企业或者外国其他组织的，符合专利法第十八条第一款的有关规定，其所属国符合专利法第十七条的有关规定。\par}

{\setlength{\parindent}{17.1bp}（5）申请人是香港、澳门或者台湾地区的个人、企业或者其他组织的，符合本指南第一部分第一章第6.1.1节的有关规定。\par}

\guidesubsection{gpos.5.3.2.2}{2.2}{不受理的情形}{\guidenote{细则 44}}{74.4bp}
专利申请有下列情形之一的，专利局不予受理：

{\setlength{\parindent}{17.1bp}（1）发明专利申请缺少请求书、说明书或者权利要求书的；实用新型专利申请缺少请求书、说明书、说明书附图或者权利要求书的；外观设计专利申请缺少请求书、图片或者照片或者简要说明的。\par}

{\setlength{\parindent}{17.1bp}（2）未使用中文的。\par}

{\setlength{\parindent}{17.1bp}（3）不符合本章第2.1节（3）中规定的受理条件的。\par}

{\setlength{\parindent}{17.1bp}（4）请求书中缺少申请人姓名或者名称，或者缺少地址的。\par}

{\setlength{\parindent}{17.1bp}\guidenote{法 17}（5）外国申请人因国籍或者居所原因，明显不具有提出专利申请的资格的。\par}

{\setlength{\parindent}{17.1bp}\guidenote{法 18.1}（6）在中国内地没有经常居所或者营业所的外国人、外国企业或者外国其他组织单独申请专利，或者作为代表人申请专利，没有委托专利代理机构的。\par}

{\setlength{\parindent}{17.1bp}（7）在中国内地没有经常居所或者营业所的香港、澳门或者台湾地区的个人、企业或者其他组织单独申请专利，或者作为代表人申请专利，没有委托专利代理机构的。\par}

{\setlength{\parindent}{17.1bp}（8）直接从外国向专利局邮寄的。\par}

{\setlength{\parindent}{17.1bp}（9）直接从香港、澳门或者台湾地区向专利局邮寄的。\par}

{\setlength{\parindent}{17.1bp}（10）专利申请类别（发明、实用新型或者外观设计）不明确或者难以确定的。\par}

{\setlength{\parindent}{17.1bp}\guidenote{细则 48.3}（11）分案申请改变原申请的类别的。\par}

\guidesubsection{gpos.5.3.2.3}{2.3}{受理与不受理程序}{}{74.4bp}
专利局受理部门收到专利申请后，应当检查和核对全部文件，作出受理或者不受理决定。

\guidesubsubsection{gpos.5.3.2.3.1}{2.3.1}{受理程序}{\guidenote{细则 43}}{74.4bp}
专利申请符合受理条件的，受理程序如下：

{\setlength{\parindent}{17.1bp}（1）确定收到日：记录受理部门收到申请文件的日期。\par}

{\setlength{\parindent}{17.1bp}（2）核实文件：清点全部文件数量，核对请求书上注明的申请文件和其他文件名称与数量，并记录核实情况。对于涉及核苷酸或者氨基酸序列的纸件发明专利申请，还应当核实是否提交了包含相应序列表的计算机可读形式的副本，例如光盘或者软盘等。通过邮寄递交的专利申请，记录文件的信件编码，并将信封存档。\par}

{\setlength{\parindent}{17.1bp}\guidenote{法 28}\guidenote[15.6bp]{细则 4.1、 4.2}（3）确定申请日：电子申请以符合要求的申请文件进入专利局指定的特定电子系统的日期为申请日。向专利局受理部门窗口直接递交的专利申请，以收到日为申请日；通过邮局邮寄递交到专利局受理部门的专利申请，以信封上的寄出邮戳日为申请日；信封上无寄出邮戳或者寄出的邮戳日不清晰或者异常的，以专利局受理部门收到日为申请日。通过速递公司递交到专利局受理部门的专利申请，以收到日为申请日。邮寄或者递交到专利局非受理部门或者个人的专利申请，以受理部门实际收到日为申请日。\par}

{\setlength{\parindent}{17.1bp}（4）给予申请号：按照专利申请的类别和专利申请的先后顺序给出相应的专利申请号。\par}

{\setlength{\parindent}{17.1bp}（5）发出通知书：发出专利申请受理通知书、缴纳申请费通知书或者收费减缴审批通知书。\par}

\guidesubsubsection{gpos.5.3.2.3.2}{2.3.2}{分案申请的受理程序}{}{102.0bp}
\guideparagraph{gpos.5.3.2.3.2.1}{2.3.2.1}{国家申请的分案申请的受理程序}{}{0bp}
对于国家申请的分案申请除按照一般专利申请的受理条件对分案申请进行受理审查外，还应当对分案申请请求书中是否填写了原申请的申请号和原申请的申请日进行审查。分案申请请求书中原申请的申请号填写正确，但未填写原申请的申请日的，以原申请号所对应的申请日为申请日。分案申请请求书中未填写原申请的申请号或者填写的原申请的申请号有误的，按照一般专利申请受理。

对符合受理条件的分案申请，专利局应当受理，给出专利申请号，以原申请的申请日为申请日，并记载分案申请递交日。

\guideparagraph{gpos.5.3.2.3.2.2}{2.3.2.2}{进入国家阶段的国际申请的分案申请的受理程序}{}{74.4bp}
国际申请进入国家阶段之后提出的分案申请，审查员除了按照一般专利申请的受理条件对分案申请进行受理审查外，还应当核实分案申请请求书中是否填写了原申请的申请日和原申请的申请号，该原申请的申请日应当是其国际申请日，原申请的申请号是进入国家阶段时专利局给予的申请号。

\guidesubsubsection{gpos.5.3.2.3.3}{2.3.3}{根据专利法实施细则第四十五条通过援引在先申请补交遗漏文件的受理程序}{}{74.4bp}
\guidenote{细则 45}根据专利法实施细则第四十五条的规定，申请人以援引在先申请的方式补交权利要求书或者说明书（实用新型附图）的，应当在首次递交专利申请时提出援引加入声明。申请人使用专利局制定的包含援引加入声明的专利请求书标准表格，即视为提出援引加入声明。在受理程序中发现发明专利申请或者实用新型专利申请缺少说明书（实用新型无附图）或者权利要求书时，如果该申请要求了优先权，专利局发出补交遗漏文件通知书。未要求优先权的，专利局发出文件不受理通知书。分案申请不适用专利法实施细则第四十五条的规定。

申请人在首次递交专利申请之日起两个月内或者收到补交遗漏文件通知书之日起两个月内提交确认援引加入声明,补交遗漏文件且满足受理条件的，专利局发出受理通知书、缴纳申请费通知书或者收费减缴审批通知书；未在规定期限内通过援引在先申请补交遗漏文件或者补交后仍不满足受理条件的，专利局发出文件不受理通知书。

\guidesubsubsection{gpos.5.3.2.3.4}{2.3.4}{不受理程序}{}{74.4bp}
专利申请不符合受理条件的，不受理程序如下：

{\setlength{\parindent}{17.1bp}（1）确定收到日：记录受理部门收到申请文件的日期。\par}

{\setlength{\parindent}{17.1bp}（2）记录不受理文档编号、不受理原因等信息并发出文件不受理通知书。\par}

{\setlength{\parindent}{17.1bp}（3）不符合受理条件的申请文件存档备查，原则上不退回当事人。\par}

在专利局受理部门窗口直接递交的专利申请，不符合受理条件的，应当直接向当事人说明原因，不予接收。

\guidesection{gpos.5.3.3}{3.}{其他文件的受理与不受理}{}{102.0bp}
\guidesubsection{gpos.5.3.3.1}{3.1}{其他文件的受理条件}{}{0bp}
申请后当事人提交的其他文件符合下列条件的，专利局应当受理：

{\setlength{\parindent}{17.1bp}（1）各文件中注明了该文件所涉及专利申请的申请号（或者专利号），并且仅涉及一件专利申请（或者专利）。\par}

{\setlength{\parindent}{17.1bp}（2）各文件用中文书写，字迹清晰、字体工整，并且用不易擦去的笔迹完成；外文证明材料附有中文清单。\par}

专利局受理部门收到申请人（或者专利权人）或者其他相关当事人递交的与专利申请有关的其他文件时，应当检查和核对全部文件。

\guidesubsection{gpos.5.3.3.2}{3.2}{其他文件的受理程序}{}{74.4bp}
其他文件符合受理条件的，受理程序如下：

{\setlength{\parindent}{17.1bp}（1）确定收到日：记录受理部门收到文件的日期。\par}

{\setlength{\parindent}{17.1bp}（2）核实文件：清点全部文件数量,并记录文件核实情况。递交文件的申请号是错号的，若依据其他信息能正确判定其正确申请号的，可以依职权予以确定；若不能予以判定的，则不予受理。\par}

{\setlength{\parindent}{17.1bp}（3）确定递交日：其他文件递交日的确定参照本章第2.3.1节第（3）项的规定。\par}

{\setlength{\parindent}{17.1bp}（4）给出收到文件回执：对于当事人在受理窗口递交的文件，经核实文件信息后，专利局出具收到文件回执，回执上应注明受理部门收到日。当事人以寄交方式递交文件的，专利局不再出具收到文件回执。\par}

\guidesubsection{gpos.5.3.3.3}{3.3}{其他文件的不受理程序}{}{74.4bp}
其他文件不符合受理条件的，按照本章第2.3.3节中规定的程序处理并发出文件不受理通知书。

\guidesection{gpos.5.3.4}{4.}{申请日的更正}{}{74.4bp}
申请人收到专利申请受理通知书之后认为该通知书上记载的申请日与邮寄该申请文件日期不一致的，可以请求专利局更正申请日。

专利局收到申请人的申请日更正请求后，应当检查更正请求是否符合下列规定：

{\setlength{\parindent}{17.1bp}（1）在递交专利申请文件之日起两个月内或者申请人收到专利申请受理通知书之日起一个月内提出。\par}

{\setlength{\parindent}{17.1bp}（2）附有收寄专利申请文件的邮局出具的寄出日期的有效证明，该证明中注明的寄出挂号号码与请求书中记录的挂号号码一致。\par}

申请日更正请求符合规定的，审查员应当发出修改更正通知书；不符合规定的，审查员应当发出视为未提出通知书，并说明理由。

当事人对专利局确定的其他文件递交日有异议的，应当提供专利局出具的收到文件回执、收寄邮局出具的证明或者其他有效证明材料。证明材料符合规定的，专利局应当重新确定递交日并修改有关数据。

\guidesection{gpos.5.3.5}{5.}{受理程序中错误的更正}{}{74.4bp}
专利局受理部门在受理工作中出现的错误一经发现，应当及时更正。

因更正影响应缴费用金额的，当事人应当按照本部分第二章第4.2.4.3节中规定的程序办理费用相关事宜。

对专利局内部错投到非受理部门的文件应当及时退回受理部门，并注明退回原因。

\guidesection{gpos.5.3.6}{6.}{查询}{}{74.4bp}
专利局受理部门收到文件后应当记录文件提交方式、提交日期、案卷编号或者挂号号码等信息。当事人除能提供专利局受理部门的收文回执或者受理通知书外，以收文登记记录为准。

查询时效为一年，自提交该文件之日起算。

\guidechapter{gpos.5.4}{第四章}{专利申请文档}{}{0bp}
\guidesection{gpos.5.4.1}{1.}{文档及组成}{}{74.4bp}
专利申请文档是在专利申请审查程序中以及专利权有效期内逐步形成、并作为原始记录保存起来以备查考的各种文件的集合，包括案卷和电子文档。专利申请文档是专利局进行审批和作出各种结论的依据。

\guidesection{gpos.5.4.2}{2.}{案卷}{}{74.4bp}
案卷包括案卷夹和案卷夹内的各种文件。

\guidesubsection{gpos.5.4.2.1}{2.1}{案卷夹}{}{74.4bp}
案卷夹用于保存文件，同时也用于记录案卷的重要内容，因此案卷夹是案卷的一个重要组成部分。

当案卷夹遭到自然或者人为的损坏需要更换新案卷夹时，应当将案卷夹上的全部记录移至新案卷夹上，并将原案卷夹与案卷同时保存，不得销毁。

\guidesubsection{gpos.5.4.2.2}{2.2}{文件}{}{74.4bp}
专利申请案卷中的文件主要来自以下几个方面：

{\setlength{\parindent}{17.1bp}（1）申请人提出专利申请时，提交的专利申请文件及其他文件。\par}

{\setlength{\parindent}{17.1bp}（2）专利局在对专利申请文件及其他文件进行审查的过程中，申请人应审查员要求作出的各种答复。\par}

{\setlength{\parindent}{17.1bp}（3）提出专利申请之后申请人主动办理各种手续时提交的文件及证明材料。\par}

{\setlength{\parindent}{17.1bp}（4）在专利申请审查程序中以及专利权有效期间内，任何人依法对专利申请（或者专利）提交的各种文件以及人民法院等部门对这些文件审理后产生的文件。\par}

{\setlength{\parindent}{17.1bp}（5）其他相关文件。\par}

上述文件经过处理、立卷、归档形成案卷的重要组成部分。

\guidesubsection{gpos.5.4.2.3}{2.3}{案卷的立卷}{}{74.4bp}
立卷时应当遵循以下原则：

{\setlength{\parindent}{17.1bp}（1）真实原则。收集的内容应当是申请人（或者专利权人）、其他相关当事人等在申请专利、专利申请的审批、授权后各个法律程序中提交的原始文件。这些文件不得替换、删除、补充和涂改。\par}

{\setlength{\parindent}{17.1bp}（2）独立原则。每一件专利申请应当建立一份独立的案卷，以该专利申请的申请号作为该案卷的案卷号，该案卷号使用于案卷存在的全过程。\par}

同一申请人（或者专利权人）对若干专利申请（或者专利）办理内容完全相同的手续时，应当分别对所有专利申请（或者专利）提出请求，这些文件将被归入各自的案卷中。申请人（或者专利权人）不得使用“参见”的方式省略文件。对于专利申请集体进行申请人（或者专利权人）名称变更或者权利转移的，证明文件副本由专利局确认后，与正本具有同等效力。

{\setlength{\parindent}{17.1bp}（3）时间顺序原则。当事人依法向专利局办理各种手续时，专利局应当对所提出的各种文件及时处理，并立卷归档。\par}

专利申请案卷应当按照各文件处理时间的先后顺序立卷。

\guidesection{gpos.5.4.3}{3.}{电子文档}{}{74.4bp}
电子文档的建立应当参照本章第2.3节的原则，并包括以下内容：

{\setlength{\parindent}{17.1bp}（1）专利局基于当事人提交的纸件文件制作的图形文件和代码化文件；\par}

{\setlength{\parindent}{17.1bp}（2）当事人按照规定形式提交的核苷酸或者氨基酸的序列表；\par}

{\setlength{\parindent}{17.1bp}（3）在专利审批程序和复审、无效宣告程序中，专利局作出的通知、决定（如补正通知书、驳回决定等）及其他文件（如发明专利申请单行本，发明专利、实用新型专利和外观设计专利单行本等）；\par}

{\setlength{\parindent}{17.1bp}（4）专利费用的相关数据；\par}

{\setlength{\parindent}{17.1bp}（5）与专利申请或者专利审批有关的法律状态及变化的历史记录；\par}

{\setlength{\parindent}{17.1bp}（6）在专利审批程序中全部著录项目及其变更的历史记录；\par}

{\setlength{\parindent}{17.1bp}（7）当事人以电子申请方式提交的电子文件；\par}

{\setlength{\parindent}{17.1bp}（8）专利权评价报告；\par}

{\setlength{\parindent}{17.1bp}（9）分类号、所属审查部门、各种标记（如优先权标记、实审请求标记、保密标记等）。\par}

\guidesection{gpos.5.4.4}{4.}{法律效力}{}{74.4bp}
专利申请文档是对专利审批、复审、无效宣告等法律程序和涉及由权利归属纠纷引起的相关程序的真实记录。

\guidesection{gpos.5.4.5}{5.}{查阅和复制}{}{102.0bp}
\guidesubsection{gpos.5.4.5.1}{5.1}{查阅和复制的原则}{}{0bp}
{\setlength{\parindent}{17.1bp}\guidenote{法 21.3}（1）专利局对公布前的发明专利申请、授权公告前的实用新型和外观设计专利申请负有保密责任。在此期间，查阅和复制请求人仅限于该案申请人及其专利代理师。\par}

{\setlength{\parindent}{17.1bp}（2）任何人均可向专利局请求查阅和复制公布后的发明专利申请案卷和授权后的实用新型和外观设计专利申请案卷。\par}

{\setlength{\parindent}{17.1bp}\guidenote{细则 145.1}（3）对于已经审结的复审案件和无效宣告案件的案卷，原则上可以查阅和复制。\par}

{\setlength{\parindent}{17.1bp}（4）专利局相关部门对尚未审结的复审和无效案卷负有保密责任。对于复审和无效宣告程序中的文件，查阅和复制请求人仅限于该案当事人。\par}

{\setlength{\parindent}{17.1bp}（5）案件结论为视为未提出、不予受理、主动撤回、视为撤回的复审和无效案卷，对于复审和无效宣告程序中的文件，查阅和复制请求人仅限于该案当事人。\par}

{\setlength{\parindent}{17.1bp}（6）专利局根据审查需要要求当事人提供的各种文件，原则上可以查阅和复制。但查阅和复制行为可能存在损害当事人合法权益，或者涉及个人隐私或者商业秘密等情形的除外。\par}

{\setlength{\parindent}{17.1bp}（7）涉及国家利益或者因专利局内部业务及管理需要在案卷中留存的有关文件，不予查阅和复制。\par}

\guidesubsection{gpos.5.4.5.2}{5.2}{允许查阅和复制的内容}{}{74.4bp}
{\setlength{\parindent}{17.1bp}（1）对于公布前的发明专利申请、授权公告前的实用新型和外观设计专利申请，该案申请人或者代理师可以查阅和复制该专利申请案卷中的有关内容，包括：申请文件，与申请直接有关的手续文件，以及在初步审查程序中向申请人发出的通知书和决定书、申请人对通知书的答复意见正文。\par}

{\setlength{\parindent}{17.1bp}（2）对于已经公布但尚未公告授予专利权的发明专利申请案卷，可以查阅和复制该专利申请案卷中的有关内容，包括：申请文件，与申请直接有关的手续文件，公布文件，在初步审查程序中向申请人发出的通知书和决定书、申请人对通知书的答复意见正文，以及在实质审查程序中向申请人发出的通知书、检索报告和决定书。\par}

{\setlength{\parindent}{17.1bp}\guidenote{细则 145.1}（3）对于已经公告授予专利权的专利申请案卷，可以查阅和复制的内容包括：申请文件，优先权文件，与申请直接有关的手续文件，发明专利申请单行本，发明专利、实用新型专利和外观设计专利单行本，专利登记簿，专利权评价报告，以及在各已审结的审查程序（包括初步审查、实质审查、复审和无效宣告等）中专利局向申请人或者有关当事人发出的通知书、检索报告和决定书、申请人或者有关当事人对通知书的答复意见。\par}

{\setlength{\parindent}{17.1bp}（4）对于处在复审程序、无效宣告程序之中尚未结案的专利申请案卷，因特殊情况需要查阅和复制的，经有关方面同意后，参照上述第（1）和第（2）项的有关规定查阅和复制专利申请案卷中进入当前审查程序以前的内容。\par}

\guidesubsection{gpos.5.4.5.3}{5.3}{查阅和复制程序}{}{74.4bp}
查阅和复制专利申请案卷中的文件，应当按照下列顺序进行：

{\setlength{\parindent}{17.1bp}（1）请求人提出书面请求并缴纳规定费用。\par}

{\setlength{\parindent}{17.1bp}（2）专利局工作人员在审核请求人出具的有关证明或者证件后，到案卷所在部门提取案卷，根据本章第5.2节的规定对案卷进行整理，取出不允许查阅和复制的文件。\par}

{\setlength{\parindent}{17.1bp}（3）与请求人约定查阅时间并发出查阅通知书。\par}

{\setlength{\parindent}{17.1bp}（4）查阅人凭查阅通知书到指定地点查阅文件，对需要复制的文件进行复制。\par}

{\setlength{\parindent}{17.1bp}（5）专利局工作人员对查阅完毕的专利申请案卷重新整理，并将请求阅档的证明原件和证件复印件存入案卷后，将该案卷退回所在部门。\par}

\guidesection{gpos.5.4.6}{6.}{案卷的保存期限和销毁}{}{102.0bp}
\guidesubsection{gpos.5.4.6.1}{6.1}{保存期限}{\guidenote{细则 145.2 及.3}}{0bp}
已结案的案卷可分成：未授权结案（视为撤回、撤回和驳回等）的案卷和授权后结案（视为放弃取得专利权、主动放弃专利权、未缴年费专利权终止、专利权期限届满和专利权被宣告全部无效等）的案卷两种。

未授权结案的案卷保存期限为二年；授权后结案的案卷保存期限为三年。保存期限自结案日起算。

有分案申请的原申请的案卷的保存期从最后结案的分案的结案日起算。

作出不受理决定的专利申请文件保存期限为一年。保存期限自不受理通知书发出之日起算。

\guidesubsection{gpos.5.4.6.2}{6.2}{销毁}{}{74.4bp}
销毁前通过计算机作出案卷销毁清册，该清册记载被销毁的案卷的案卷号、基本著录项目、销毁日期。清册经主管局长签署同意销毁后，由主管案卷部门实施销毁工作。

\guidechapter{gpos.5.5}{第五章}{保密申请与向外国申请专利的保密审查}{}{0bp}
\guidesection{gpos.5.5.1}{1.}{保密的范围}{}{74.4bp}
专利法第四条规定的保密范围是涉及国家安全或者重大利益两个方面的发明创造。

根据专利法实施细则第七条第一款的规定，专利局受理的专利申请涉及国防利益需要保密的，应当及时移交国防知识产权局进行审查。

根据专利法实施细则第七条第二款的规定，专利局认为其受理的发明或者实用新型专利申请涉及国防利益以外的国家安全或者重大利益需要保密的，应当及时作出按照保密专利申请处理的决定，并通知申请人。

\guidesection{gpos.5.5.2}{2.}{保密的基准}{}{74.4bp}
保密的基准按照国家有关规定执行。

\guidesection{gpos.5.5.3}{3.}{专利申请的保密确定}{}{129.6bp}
\guidesubsection{gpos.5.5.3.1}{3.1}{申请人提出保密请求的保密确定}{}{0bp}
\guidesubsubsection{gpos.5.5.3.1.1}{3.1.1}{保密请求的提出}{}{0bp}
申请人认为其发明或者实用新型专利申请涉及除国防利益以外的国家安全或者重大利益需要保密的，应当在提出专利申请的同时，在请求书上提出保密请求，其申请文件应当以纸件形式提交。申请人也可以在发明专利申请作好公布准备之前，或者实用新型专利申请作好授权公告准备之前，提出保密请求。

申请人在提出保密请求之前已确定其申请的内容涉及国家安全或者重大利益需要保密的，应当提交有定密权限的机关、单位出具的保密证明材料。

保密证明材料中应当写明发明创造名称、申请人姓名或者名称、密级、保密期限、保密原因及保密要点、定密责任人、定密机关或者单位的联系人和联系电话、申请人机要通信地址以及定密日期。该文件还应当加盖定密机关或者单位的公章。

\guidesubsubsection{gpos.5.5.3.1.2}{3.1.2}{保密的确定}{}{74.4bp}
审查员应当根据保密基准对专利申请进行审查，并根据不同情况确定是否需要保密。

{\setlength{\parindent}{17.1bp}\guidenote{细则 7.1}（1）专利申请的内容涉及国防利益的，由国防知识产权局进行保密确定。需要保密的，应当及时移交国防知识产权局进行审查，审查员向申请人发出专利申请移交国防知识产权局通知书；不需要保密的，审查员应当发出保密审批通知书，通知申请人该专利申请不予保密，按照一般专利申请处理。\par}

{\setlength{\parindent}{17.1bp}\guidenote{细则 7.2}（2）发明或者实用新型内容涉及国防利益以外的国家安全或者重大利益的，由专利局进行保密确定，必要时可以邀请相关领域的技术专家协助确定。审查员根据保密确定的结果发出保密审批通知书，需要保密的，通知申请人该专利申请予以保密，按照保密专利申请处理；不需要保密的，通知申请人该专利申请不予保密，按照一般专利申请处理。\par}

\guidesubsection{gpos.5.5.3.2}{3.2}{专利局自行进行的保密确定}{}{74.4bp}
分类审查员在对发明或者实用新型专利申请进行分类时，应当将发明内容可能涉及国家安全或者重大利益，但申请人未提出保密请求的发明或者实用新型专利申请挑选出来。审查员应当参照本章第3.1.2节的规定，对上述专利申请进行保密确定。

对于已确定为保密专利申请的电子申请，审查员应当将该专利申请转为纸件形式继续审查并通知申请人。申请人此后应当以纸件形式向专利局或者国防知识产权局递交各种文件。

\guidesection{gpos.5.5.4}{4.}{保密专利申请的审批流程}{\guidenote{细则 7}}{74.4bp}
{\setlength{\parindent}{17.1bp}\guidenote{细则 61}（1）涉及国防利益需要保密的专利申请，由国防知识产权局进行审查，经审查没有发现驳回理由的，由专利局根据国防知识产权局的审查意见作出授予国防专利权的决定，并委托国防知识产权局颁发国防专利证书，同时在专利公报上公告国防专利的专利号、申请日和授权公告日。\par}

国防知识产权局专利复审委员会作出宣告国防专利权无效决定的，专利局应当在专利公报上公告专利号、授权公告日、无效宣告决定号和无效宣告决定日。

{\setlength{\parindent}{17.1bp}（2）涉及国防利益以外的国家安全或者重大利益需要保密的发明或者实用新型专利申请，由专利局按照以下程序进行审查和管理。\par}

审查员应当对确定需要保密的专利申请案卷作出保密标记，在对该专利申请作出解密决定之前，对其进行保密管理。

保密专利申请的初步审查和实质审查均由专利局指定的审查员进行。

对于发明专利申请，初步审查和实质审查按照与一般发明专利申请相同的基准进行。初步审查合格的保密专利申请不予公布，实质审查请求符合规定的，直接进入实质审查程序。经实质审查没有发现驳回理由的，作出授予保密发明专利权的决定，并发出授予发明专利权通知书和办理登记手续通知书。

对于实用新型专利申请，初步审查按照与一般实用新型专利申请相同的基准进行。经初步审查没有发现驳回理由的，作出授予保密实用新型专利权的决定，并发出授予实用新型专利权通知书和办理登记手续通知书。

保密专利申请的授权公告仅公布专利号、申请日和授权公告日。

\guidesection{gpos.5.5.5}{5.}{专利申请（或者专利）的解密程序}{}{102.0bp}
\guidesubsection{gpos.5.5.5.1}{5.1}{申请人（或者专利权人）提出解密请求}{}{0bp}
保密专利申请的申请人或者保密专利的专利权人可以书面提出解密请求。提出保密请求时提交了有关部门确定密级的相关文件的，申请人（或者专利权人）提出解密请求时，应当附具原确定密级的机关、单位同意解密的证明文件。

专利局对提出解密请求的保密专利申请（或者专利）进行解密确定，并将结果通知申请人。

\guidesubsection{gpos.5.5.5.2}{5.2}{专利局定期解密}{}{74.4bp}
专利局每两年对保密专利申请（或者专利）进行一次复查，经复查认为不需要继续保密的，通知申请人予以解密。

\guidesubsection{gpos.5.5.5.3}{5.3}{解密后的处理}{}{74.4bp}
审查员应当对已经解密的专利申请（或者专利）作出解密标记。发明专利申请解密后，尚未被授予专利权的，按照一般发明专利申请进行审查和管理，符合公布条件的，应当予以公布，并出版发明专利申请单行本；实用新型专利申请解密后，尚未被授予专利权的，按照一般实用新型专利申请进行审查和管理。

发明或者实用新型专利解密后，应当进行解密公告、出版发明或者实用新型专利单行本，并按照一般专利进行管理。

\guidesubsection{gpos.5.5.5.4}{5.4}{解密国防专利的接收和处理}{}{74.4bp}
专利局接收国防知识产权局移交的解密国防专利后，应当及时对该专利进行解密公告、出版发明专利单行本，并将该专利转为一般专利进行管理。

\guidesection{gpos.5.5.6}{6.}{向外国申请专利的保密审查}{}{74.4bp}
\guidenote{法 19.1}专利法第十九条第一款规定，任何单位或者个人将在中国完成的发明或者实用新型向外国申请专利的，应当事先报经专利局进行保密审查。

\guidenote{法 19.4}专利法第十九条第四款规定，对违反本条第一款规定向外国申请专利的发明或者实用新型，在中国申请专利的，不授予专利权。

\guidenote{细则 8}根据专利法实施细则第八条的规定，任何单位或者个人将在中国完成的发明或者实用新型向外国申请专利的，应当采用下列方式之一请求专利局进行保密审查：

{\setlength{\parindent}{17.1bp}（1）直接向外国申请专利或者向有关国外机构提交专利国际申请的，应当事先向专利局提出请求，并详细说明其技术方案；\par}

{\setlength{\parindent}{17.1bp}（2）向专利局申请专利后拟向外国申请专利或者向有关国外机构提交专利国际申请的，应当在向外国申请专利或者向有关国外机构提交专利国际申请前向专利局提出请求。\par}

向专利局提交专利国际申请的，视为同时提出了保密审查请求。

上述规定中所述的向外国申请专利是指向外国国家或外国政府间专利合作组织设立的专利主管机构提交专利申请，向有关国外机构提交专利国际申请是指向作为PCT受理局的外国国家或外国政府间专利合作组织设立的专利主管机构或世界知识产权组织国际局提交专利国际申请。

\guidesubsection{gpos.5.5.6.1}{6.1}{准备直接向外国申请专利的保密审查}{}{102.0bp}
\guidesubsubsection{gpos.5.5.6.1.1}{6.1.1}{保密审查请求的提出}{}{0bp}
\guidenote{细则 8.2}\guidenote[15.6bp]{细则 20}向外国申请专利保密审查请求的文件应当包括向外国申请专利保密审查请求书和技术方案说明书。请求书和技术方案说明书应当使用中文，请求人可以同时提交相应的外文文本供审查员参考。技术方案说明书应当与向外国申请专利的内容一致。技术方案说明书可以参照专利法实施细则第二十条的规定撰写，并符合本部分第一章的其他规定。

\guidesubsubsection{gpos.5.5.6.1.2}{6.1.2}{保密审查}{}{74.4bp}
审查员对向外国申请专利保密审查请求文件进行初步保密审查。请求文件形式不符合规定的，审查员应当在请求递交日起两个月内通知请求人该向外国申请专利保密审查请求视为未提出，情况复杂的，可以在请求递交日起四个月内通知请求人，请求人可以重新提出符合规定的向外国申请专利保密审查请求。

技术方案明显不需要保密的，审查员应当在请求递交日起两个月内通知请求人可以就该技术方案向外国申请专利，情况复杂的，可以在请求递交日起四个月内通知请求人。技术方案可能需要保密的，审查员应当将需作进一步保密审查、暂缓向外国申请专利的审查意见通知请求人。审查员应当在请求递交日起两个月内发出向外国申请专利保密审查意见通知书，情况复杂的，在请求递交日起四个月内发出通知书，将上述审查结论通知请求人。

已通知请求人暂缓向外国申请专利的，审查员应当作进一步保密审查，必要时可以邀请相关领域的技术专家协助审查。审查员根据保密审查的结论在请求递交日起四个月内发出向外国申请专利保密审查决定，情况复杂的，在请求递交日起六个月内发出决定，将是否同意就该技术方案向外国申请专利的审查结果通知请求人。

\guidesubsection{gpos.5.5.6.2}{6.2}{申请专利后拟向外国申请专利的保密审查}{}{102.0bp}
\guidesubsubsection{gpos.5.5.6.2.1}{6.2.1}{保密审查请求的提出}{}{0bp}
\guidenote{细则 8.2}申请人拟在向专利局申请专利后又向外国申请专利的，应当在提交专利申请同时或者之后提交向外国申请专利保密审查请求书。未按上述规定提出请求的，视为未提出请求。向外国申请专利的内容应当与该专利申请的内容一致。

\guidesubsubsection{gpos.5.5.6.2.2}{6.2.2}{保密审查}{}{74.4bp}
对提出向外国申请专利保密审查请求的专利申请，审查员应当参见本章第6.1.2节中的规定进行保密审查。

\guidesubsection{gpos.5.5.6.3}{6.3}{国际申请的保密审查}{}{102.0bp}
\guidesubsubsection{gpos.5.5.6.3.1}{6.3.1}{保密审查请求的提出}{\guidenote{细则 8.2}}{0bp}
申请人向专利局提交国际申请的，视为同时提出向外国申请专利保密审查请求。

\guidesubsubsection{gpos.5.5.6.3.2}{6.3.2}{保密审查}{}{74.4bp}
国际申请不需要保密的，审查员应当按照正常国际阶段程序进行处理。国际申请需要保密的，审查员应当自申请日起三个月内发出因国家安全原因不再传送登记本和检索本的通知书，通知申请人和国际局该申请将不再作为国际申请处理，终止国际阶段程序。申请人收到上述通知的，不得就该申请的内容向外国申请专利。

\guidechapter{gpos.5.6}{第六章}{通知和决定}{}{0bp}
\guidesection{gpos.5.6.1}{1.}{通知和决定的产生}{}{102.0bp}
\guidesubsection{gpos.5.6.1.1}{1.1}{通知和决定的种类}{}{0bp}
在专利申请的审批程序、复审程序、无效宣告程序以及专利法及其实施细则规定的其他程序中，审查员根据不同情况，将作出各种通知和决定。这些通知和决定主要包括：专利申请受理通知书、审查意见通知书、补正通知书、手续合格通知书、视为撤回通知书、恢复权利请求审批通知书、发明专利申请实质审查请求期限届满前通知书、缴费通知书、收费减缴审批通知书、发明专利申请初步审查合格通知书、发明专利申请公布通知书、发明专利申请进入实质审查阶段通知书、授予发明专利权通知书、授予实用新型专利权通知书、授予外观设计专利权通知书、办理登记手续通知书、视为放弃取得专利权通知书、专利权终止通知书、驳回决定、复审决定书、无效宣告请求审查决定等。

\guidesubsection{gpos.5.6.1.2}{1.2}{通知和决定的撰写}{}{74.4bp}
撰写通知和决定应当符合专利法及其实施细则和本指南的有关规定。

除本指南中其他章节作出专门规定之外，通知和决定一般应当包括：收件人信息、著录项目、通知或者决定的内容、署名和盖章、发文日期。其中：

{\setlength{\parindent}{17.1bp}（1）收件人信息包括：收件人地址、邮政编码、收件人姓名。\par}

{\setlength{\parindent}{17.1bp}（2）著录项目包括：申请号（或者专利号）、发明创造名称、全体申请人（或者专利权人）姓名或者名称，如果是无效宣告、中止程序中的通知书，还应当包括全体请求人的姓名或者名称。\par}

{\setlength{\parindent}{17.1bp}（3）通知、决定的内容包括：通知或者决定的名称及正文。在作出不利于当事人的通知或者决定时，应当说明理由，必要时指明后续法律程序。\par}

{\setlength{\parindent}{17.1bp}（4）署名和盖章：通知和决定应当有审查员署名；需要审核的，还应当有审核人员署名；发出的通知和决定均应当加盖国家知识产权局相应的业务用章。\par}

\guidesection{gpos.5.6.2}{2.}{通知和决定的送达}{}{129.6bp}
\guidesubsection{gpos.5.6.2.1}{2.1}{送达方式}{\guidenote{细则 4}}{0bp}
\guidesubsubsection{gpos.5.6.2.1.1}{2.1.1}{邮寄}{}{0bp}
邮寄送达文件是指通过邮局把通知和决定送交当事人。除另有规定外，邮寄的文件应当挂号，并应当在计算机中登记挂号的号码、收件人地址和姓名、文件类别、所涉及的专利申请号、发文日期、发文部门。邮寄被退回的函件要登记退函日期。

\guidesubsubsection{gpos.5.6.2.1.2}{2.1.2}{直接送交}{}{74.4bp}
经专利局同意，专利代理机构可以在专利局指定的时间和地点，按时接收通知和决定。特殊情况下经专利局同意，当事人本人也可以在专利局指定的时间和地点接收通知和决定。

除受理窗口当面交付受理通知书和文件回执外，当面交付其他文件时应当办理登记签收手续。特殊情况下，应当由当事人在申请案卷上签字或者盖章，并记录当事人身份证件的名称、号码和签发单位。

\guidesubsubsection{gpos.5.6.2.1.3}{2.1.3}{电子形式}{}{74.4bp}
对于以电子形式提交的专利申请，专利局通过指定的特定电子系统以电子形式向申请人送达各种通知书、决定和其他文件。

\guidesubsubsection{gpos.5.6.2.1.4}{2.1.4}{公告送达}{}{74.4bp}
专利局发出的通知和决定被退回的，审查员应当与文档核对；如果确定文件因送交地址不清或者存在其他原因无法再次邮寄的，应当在专利公报上通过公告方式通知当事人。自公告之日起满一个月，该文件视为已经送达。

\guidesubsection{gpos.5.6.2.2}{2.2}{收件人}{}{102.0bp}
\guidesubsubsection{gpos.5.6.2.2.1}{2.2.1}{当事人未委托专利代理机构}{}{0bp}
当事人未委托专利代理机构的，通知和决定的收件人为请求书中填写的联系人。若请求书中未填写联系人的，收件人为当事人；当事人有两个以上时，请求书中另有声明指定非第一署名当事人为代表人的，收件人为该代表人；除此之外，收件人为请求书中第一署名当事人。

\guidesubsubsection{gpos.5.6.2.2.2}{2.2.2}{当事人已委托专利代理机构}{}{74.4bp}
当事人委托了专利代理机构的，通知和决定的收件人为该专利代理机构指定的专利代理师。专利代理师有两个的，收件人为该两名专利代理师。

\guidesubsubsection{gpos.5.6.2.2.3}{2.2.3}{其他情况}{}{74.4bp}
当事人无民事行为能力的，在专利局已被告知的情况下，通知和决定的收件人是法定监护人或者法定代理人。

\guidesubsection{gpos.5.6.2.3}{2.3}{送达日}{\guidenote{细则 4}}{102.0bp}
\guidesubsubsection{gpos.5.6.2.3.1}{2.3.1}{邮寄、直接送交和电子形式送达}{}{0bp}
通过邮寄、直接送交的通知和决定，自发文日起满十五日推定为当事人收到通知和决定之日。对于通过邮寄的通知和决定，当事人提供证据能够证明实际收到文件的日期的，以实际收到日为准。

通过电子形式送达的通知和决定，以进入当事人认可的电子系统的日期为送达日。进入当事人认可的电子系统的日期与通知书和决定的发文日不一致时，除当事人能提供证据外，该通知书和决定的发文日推定为送达日。

\guidesubsubsection{gpos.5.6.2.3.2}{2.3.2}{公告送达}{}{74.4bp}
通知和决定是通过在专利公报上公告方式通知当事人的，以公告之日起满一个月推定为送达日。当事人见到公告后可以向专利局提供详细地址，要求重新邮寄有关文件，但仍以自公告之日起满一个月为送达日。

\guidesection{gpos.5.6.3}{3.}{退件的处理和文件的查询}{}{102.0bp}
\guidesubsection{gpos.5.6.3.1}{3.1}{退件的处理}{}{0bp}
邮寄退回的通知和决定由发文部门作计算机登录，再转送相关部门进行处理。

处理退件首先应当根据申请文档中申请人、专利代理机构提供的各种文件进行分析，查清退件的原因。能够重新确定正确地址和收件人的，更正后重新发出。

退件经过处理仍无法邮寄或者再次被退回时，根据通知和决定的性质，必要时采用公告的方式送达当事人。

退件（连同信封）应当存档。

\guidesubsection{gpos.5.6.3.2}{3.2}{文件的查询}{}{74.4bp}
当事人陈述未收到专利局的某一通知和决定的，应当由退件处理部门进行查询。查询首先在专利局发文部门进行，查询结果（包括通知和决定的发文日期、挂号号码和收件人）应当由退件处理部门通知当事人。

当事人需要进一步了解送达情况的，应当办理邮路查询手续，由发文部门通过当地邮局查询收件人所在邮政部门。查询结果表明未送达的责任在专利局或者邮局的，应当按照新的发文日重新发出有关通知和决定；查询结果表明未送达的责任在收件人所在单位收发部门或者收件人本人及其有关人员的，专利局可以根据当事人的请求重新发出有关通知和决定的复印件，但不得变更发文日。

邮路查询时效为一年，自发文日起计算。

\guidechapter{gpos.5.7}{第七章}{期限、权利的恢复、中止、审查的顺序}{}{0bp}
\guidesection{gpos.5.7.1}{1.}{期限的种类}{}{102.0bp}
\guidesubsection{gpos.5.7.1.1}{1.1}{法定期限}{}{0bp}
法定期限是指专利法及其实施细则规定的各种期限。例如，发明专利申请的实质审查请求期限（专利法第三十五条第一款的规定）、申请人办理登记手续的期限（专利法实施细则第六十条第一款的规定）。

\guidesubsection{gpos.5.7.1.2}{1.2}{指定期限}{}{74.4bp}
指定期限是指审查员在根据专利法及其实施细则作出的各种通知中，规定申请人（或专利权人）、其他当事人作出答复或者进行某种行为的期限。例如，根据专利法第三十七条的规定，专利局对发明专利申请进行实质审查后，认为不符合专利法规定的，应当通知申请人，要求其在指定的期限内陈述意见，或者对其申请进行修改，该期限由审查员指定。又如，根据专利法实施细则第三条第二款的规定，当事人根据专利法及其实施细则规定提交的各种证件和证明文件是外文的，专利局认为必要时，可以要求当事人在指定期限内提交中文译文，该期限也由审查员指定。

指定期限一般为两个月。发明专利申请的实质审查程序中，申请人答复第一次审查意见通知书的期限为四个月。对于外观设计国际申请，申请人答复驳回通知的期限为四个月。对于较为简单的行为，也可以给予一个月或更短的期限。

\guidesection{gpos.5.7.2}{2.}{期限的计算}{}{102.0bp}
\guidesubsection{gpos.5.7.2.1}{2.1}{期限的起算日}{}{0bp}
{\setlength{\parindent}{17.1bp}（1）自申请日、优先权日、授权公告日等固定日期起计算\par}

大部分法定期限是自申请日、优先权日、授权公告日等固定日期起计算的。例如，专利法第四十二条规定的专利权的期限均自申请日起计算。专利法第二十九条第一款规定要求外国优先权的发明或者实用新型专利申请应当在十二个月内提出，该期限的起算日为在外国第一次提出专利申请之日（优先权日）。

{\setlength{\parindent}{17.1bp}（2）自通知和决定的送达日起计算\par}

全部指定期限和部分法定期限自通知和决定的送达日起计算。送达日的相关规定参见本部分第六章第2.3节。例如，审查员根据专利法第三十七条的规定指定申请人陈述意见或者修改其申请的期限（指定期限）是自审查意见通知书送达申请人之日起计算；专利法实施细则第六十条第一款规定的申请人办理登记手续的期限（法定期限）是自授予专利权通知书送达申请人之日起计算。

\guidesubsection{gpos.5.7.2.2}{2.2}{期限的届满日}{}{74.4bp}
期限起算日加上法定或者指定的期限即为期限的届满日。相应的行为应当在期限届满日之前、最迟在届满日当天完成。

\guidesubsection{gpos.5.7.2.3}{2.3}{期限的计算}{\guidenote{细则 5}}{74.4bp}
期限开始的当日不计算在期限内，自下一日开始计算。期限以年或者月计算的，以其最后一月的相应日（与起算日相对应的日期）为期限届满日；该月无相应日的，以该月最后一日为期限届满日。例如，一件发明专利申请的申请日为1998年6月1日，其实质审查请求期限的届满日应当是2001年6月1日。又如，专利局于2008年6月6日通过邮寄发出审查意见通知书，指定期限两个月，其推定送达日是2008年6月21日（遇休假日不顺延），则期限届满日应当是2008年8月21日。再如，专利局通过指定的特定电子系统于2022年12月31日发出的通知书，其推定送达日是2022年12月31日，如果该通知书的指定期限为两个月，则期限届满日应当是2023年2月28日。

期限届满日是法定休假日或者移用周休息日的，以法定休假日或者移用周休息日后的第一个工作日为期限届满日，该第一个工作日为周休息日的，期限届满日顺延至星期一。法定休假日包括国务院发布的《全国年节及纪念日放假办法》第二条规定的全体公民放假的节日和《国务院关于职工工作时间的规定》第七条第一款规定的周休息日。

\guidesection{gpos.5.7.3}{3.}{期限的监视}{}{102.0bp}
\guidesubsection{gpos.5.7.3.1}{3.1}{期限的确定}{}{0bp}
各种期限均自期限起算日确定。例如，申请人提出专利申请，并确定了其申请日后，在建立专利申请文档的同时确定自申请日起算的各种期限；审查员在作出各种与期限有关的通知和决定时，确定自该通知和决定送达日起算的答复期限。

\guidesubsection{gpos.5.7.3.2}{3.2}{期限监视方式}{}{74.4bp}
各种期限的监视一般由计算机系统进行。申请人办理与期限有关的手续后，在计算机系统中应当记录办理手续的日期，并将该日期与期限届满日进行比较，确定该手续在期限方面的合法性。

期限以日为单位进行监视并及时处理。期限届满日起满一个月尚未销去的期限，应当予以处理，作出相应处理决定。例如，专利局于2001年9月4日通过邮寄发出补正通知书指定申请人于一个月内提交优先权转让证明文件的中文译本，该通知书的推定送达日为2001年9月19日，期限届满日为2001年10 月19 日，如果专利局一直未收到申请人提交的中文译本，应当于2001年11月19日后针对该期限进行处理，并发出视为未要求优先权通知书。

\guidesubsection{gpos.5.7.3.3}{3.3}{期限届满的通知}{}{74.4bp}
{\setlength{\parindent}{17.1bp}（1）发明专利申请实质审查请求期限届满前三个月，对尚未提出实质审查请求或者尚未缴纳实质审查费的发明专利申请发出发明专利申请实质审查请求期限届满前通知书，通知申请人办理有关手续。\par}

{\setlength{\parindent}{17.1bp}（2）专利年费缴纳期限届满后一个月，对尚未缴纳相关费用的专利发出缴费通知书，通知专利权人在专利法实施细则第一百一十五条规定的滞纳期内缴纳相关费用及滞纳金。\par}

{\setlength{\parindent}{17.1bp}（3）其他期限届满前不发出通知书提示。\par}

\guidesection{gpos.5.7.4}{4.}{期限的延长}{}{102.0bp}
\guidesubsection{gpos.5.7.4.1}{4.1}{延长期限请求}{}{0bp}
\guidenote{细则 6.4 及 75}当事人因正当理由不能在期限内进行或者完成某一行为或者程序时，可以请求延长期限。可以请求延长的期限仅限于指定期限。但在无效宣告程序中，复审和无效审理部门指定的期限不得延长。

\guidenote{细则 6.4及 116.2}请求延长期限的，应当在期限届满前提交延长期限请求书，说明理由，并缴纳延长期限请求费。延长期限请求费以月计算。

\guidesubsection{gpos.5.7.4.2}{4.2}{延长期限请求的批准}{}{74.4bp}
延长期限请求由作出相应通知和决定的部门或者流程管理部门进行审批。

延长的期限不足一个月的，以一个月计算。延长的期限不得超过两个月。对同一通知或者决定中指定的期限一般只允许延长一次。

延长期限请求不符合规定的，审查员应当发出延长期限审批通知书，并说明不予延长期限的理由；符合规定的，审查员应当发出延长期限审批通知书，在计算机系统中更改该期限的届满日，继续监视该期限。

国家出现紧急状态或者非常情况时，专利局可以依法延长专利法实施细则规定的期限和专利局指定的期限或者简化当事人应当履行的相关手续，并以公告等形式告知当事人。

\guidesection{gpos.5.7.5}{5.}{耽误期限的处置}{}{102.0bp}
\guidesubsection{gpos.5.7.5.1}{5.1}{作出处分决定前的审核}{}{0bp}
申请人（或专利权人）耽误期限的后果是丧失各种相应的权利，这些权利主要包括：专利申请权（或专利权）、优先权等。

审查员在作出各种处分决定前，应当对是否需要作出该决定进行复核，当确认申请人（或专利权人）在规定期限之内未完成应当完成的行为时，再作出相应的处分决定。

\guidesubsection{gpos.5.7.5.2}{5.2}{处分决定}{}{74.4bp}
因耽误期限作出的处分决定主要包括：视为撤回专利申请权、视为放弃取得专利权的权利、专利权终止、不予受理、视为未提出请求和视为未要求优先权等。

处分决定的撰写应当符合本部分第六章第1.2节的规定，并自期限届满日起满一个月后作出。

\guidesubsection{gpos.5.7.5.3}{5.3}{作出处分决定后的处理}{}{74.4bp}
处分决定不影响专利申请权（或专利权）的，原程序继续进行。

处分决定作出后，专利申请权（或专利权）丧失的，应当按照规定给予两个月（自该处分决定的送达日起算）的恢复权利请求期限，期满未提出恢复权利请求或者恢复权利请求不符合规定的，自处分通知书发出之日起四个月（涉及复审或者无效宣告程序的为六个月）后分别按照以下情形处理：

{\setlength{\parindent}{17.1bp}（1）处分决定涉及尚未公开的专利申请的，应当对处分决定再次复核，确定无误的，将专利申请进行失效处理。\par}

{\setlength{\parindent}{17.1bp}（2）处分决定涉及已公布的发明专利申请或者已公告的专利的，应当对处分决定再次复核，确定无误的，在专利公报上公告相应处分决定，将专利申请（或专利）进行失效处理。\par}

作出丧失专利申请权（或专利权）的处分决定后又收到有关文件表明相关手续已在规定的期限内完成的，流程部门应当及时撤销有关处分决定，发出修改更正通知书，处分决定已公告的还应当作出公告更正。

\guidesection{gpos.5.7.6}{6.}{权利的恢复}{\guidenote{细则 6}}{102.0bp}
\guidesubsection{gpos.5.7.6.1}{6.1}{适用范围}{}{0bp}
专利法实施细则第六条第一款和第二款规定了当事人因耽误期限而丧失权利之后，请求恢复其权利的条件。该条第五款又规定，不丧失新颖性的宽限期、优先权期限、专利权期限和侵权诉讼时效这四种期限被耽误而造成的权利丧失，不能请求恢复权利。

\guidesubsection{gpos.5.7.6.2}{6.2}{手续}{}{74.4bp}
根据专利法实施细则第六条第二款规定请求恢复权利的，应当自收到专利局的处分决定之日起两个月内、或者请求复审的期限届满之日起两个月内提交恢复权利请求书，说明理由，并同时缴纳恢复权利请求费；根据专利法实施细则第六条第一款规定请求恢复权利的，应当自障碍消除之日起两个月内且自期限届满之日起两年内提交恢复权利请求书，说明理由，必要时还应当附具有关证明文件。

当事人在请求恢复权利的同时，应当办理权利丧失前应当办理的相应手续，消除造成权利丧失的原因。例如，申请人因未缴纳申请费，其专利申请被视为撤回后，在请求恢复其申请权的同时，还应当补缴规定的申请费。

\guidesubsection{gpos.5.7.6.3}{6.3}{审批}{}{74.4bp}
审查员应当按照本章第6.1 节和第6.2 节的规定对恢复权利的请求进行审查。

{\setlength{\parindent}{17.1bp}（1）恢复权利的请求符合规定的，应当准予恢复权利，并发出恢复权利请求审批通知书。申请人提交信函表明请求恢复权利的意愿，只要写明申请号（或专利号）并且签字或者盖章符合要求的，可视为合格的恢复权利请求书。\par}

{\setlength{\parindent}{17.1bp}（2）已在规定期限内提交了书面请求或缴足恢复权利请求费，但仍不符合规定的，审查员应当发出办理恢复权利手续补正通知书，要求当事人在指定期限之内补正或者补办有关手续，补正或者补办的手续符合规定的，应当准予恢复权利，并发出恢复权利请求审批通知书。期满未补正或者经补正仍不符合规定的，不予恢复，发出恢复权利请求审批通知书，并说明不予恢复的理由。\par}

经专利局同意恢复专利申请权（或专利权）的，继续专利审批程序。对于已公告过处分决定的，还应当在专利公报上公告恢复权利的决定。

\guidesection{gpos.5.7.7}{7.}{中止程序}{\guidenote{细则 103.1 及}}{74.4bp}
\guidenote{104}中止，是指当地方知识产权管理部门或者人民法院受理了专利申请权（或专利权）权属纠纷，或者人民法院裁定对专利申请权（或专利权）采取财产保全措施时，专利局根据权属纠纷的当事人的请求或者人民法院的要求中止有关程序的行为。

\guidesubsection{gpos.5.7.7.1}{7.1}{请求中止的条件}{}{74.4bp}
请求专利局中止有关程序应当符合下列条件：

{\setlength{\parindent}{17.1bp}\guidenote{细则 103.1}（1）当事人请求中止的，专利申请权（或专利权）权属纠纷已被地方知识产权管理部门或者人民法院受理；人民法院要求协助执行对专利申请权（或专利权）采取财产保全措施的，应当已作出财产保全的民事裁定。\par}

{\setlength{\parindent}{17.1bp}\guidenote{细则 104}（2）中止的请求人是权属纠纷的当事人或者对专利申请权（或专利权）采取财产保全措施的人民法院。\par}

\guidesubsection{gpos.5.7.7.2}{7.2}{中止的范围}{}{74.4bp}
\guidenote{细则 105}中止的范围是指：

{\setlength{\parindent}{17.1bp}（1）暂停专利申请的初步审查、实质审查、复审、授予专利权程序和专利权无效宣告程序；\par}

{\setlength{\parindent}{17.1bp}（2）暂停视为撤回专利申请、视为放弃取得专利权、未缴年费终止专利权等程序；\par}

{\setlength{\parindent}{17.1bp}（3）暂停办理撤回专利申请、放弃专利权、变更申请人（或专利权人）的姓名或者名称、转移专利申请权（或专利权）、专利权质押登记等手续。\par}

中止请求批准前已进入公布或者公告准备的，该程序不受中止的影响。

\guidesubsection{gpos.5.7.7.3}{7.3}{请求中止的手续和审批}{}{129.6bp}
\guidesubsubsection{gpos.5.7.7.3.1}{7.3.1}{权属纠纷的当事人请求的中止}{}{0bp}
\guideparagraph{gpos.5.7.7.3.1.1}{7.3.1.1}{权属纠纷的当事人请求中止的手续}{}{0bp}
\guidenote{细则 103.2}专利申请权（或专利权）权属纠纷的当事人请求专利局中止有关程序的，应当符合下列规定：

{\setlength{\parindent}{17.1bp}（1）提交中止程序请求书并说明理由；\par}

{\setlength{\parindent}{17.1bp}（2）附具证明文件，即地方知识产权管理部门或者人民法院的写明专利申请号（或专利号）的有关受理文件正本或者副本。\par}

\guideparagraph{gpos.5.7.7.3.1.2}{7.3.1.2}{权属纠纷的当事人请求中止的审批及处理}{}{74.4bp}
专利局收到当事人提出的中止程序请求书和有关证明后，专利局的流程管理部门应当审查是否满足下列各项条件：

{\setlength{\parindent}{17.1bp}（1）请求中止的专利申请（或专利）未丧失权利，涉及无效宣告程序的除外；\par}

{\setlength{\parindent}{17.1bp}（2）未执行中止程序；\par}

{\setlength{\parindent}{17.1bp}（3）请求是由有关证明文件中所记载的权属纠纷当事人提出；\par}

{\setlength{\parindent}{17.1bp}（4）受理权属纠纷的机关对该专利申请（或专利）权属纠纷案有管辖权；\par}

{\setlength{\parindent}{17.1bp}（5）证明文件中记载的申请号（或专利号）、发明创造名称和权利人与请求中止的专利申请（或专利）记载的内容一致；\par}

{\setlength{\parindent}{17.1bp}（6）中止请求书与证明文件其他方面符合规定的形式要求。\par}

不满足上述第（1）至（5）项条件的，审查员应当向中止程序请求人发出视为未提出通知书。不满足上述第（6）项条件的，例如中止程序请求书不符合格式要求或者提交的证明文件不是正本或者副本的，审查员应当发出办理手续补正通知书，通知中止程序请求人在一个月的期限内补正其缺陷。补正期限内，暂停有关程序。期满未补正的或者补正后仍未能消除缺陷的，应当向中止程序请求人发出视为未提出通知书，恢复有关程序。

满足上述条件或者经补正后满足上述条件的，审查员向专利申请（或专利）权属纠纷的双方当事人发出中止程序请求审批通知书，并告知中止期限的起止日期（自提出中止请求之日起）。

但是，对于处于无效宣告程序中的专利，专利局的流程管理部门完成形式审查后，由复审和无效审理部门作进一步审查。具备下列情形之一的，可以不中止专利权无效宣告程序：

{\setlength{\parindent}{17.1bp}（1）根据已进行的审查工作能够作出无效宣告审查决定的；\par}

{\setlength{\parindent}{17.1bp}（2）权属纠纷当事人依据的理由明显不充分，也未提交足以证明确有权属纠纷存在的证据的；\par}

{\setlength{\parindent}{17.1bp}（3）有证据表明，中止专利权无效宣告程序将明显损害当事人利益或者公共利益的；\par}

{\setlength{\parindent}{17.1bp}（4）有证据表明，中止程序的请求明显具有不诚信、不正当行为的。\par}

\guidesubsubsection{gpos.5.7.7.3.2}{7.3.2}{因人民法院要求协助执行财产保全的中止}{}{102.0bp}
\guideparagraph{gpos.5.7.7.3.2.1}{7.3.2.1}{因协助执行财产保全而中止的手续}{\guidenote{细则 104}}{0bp}
因人民法院要求协助执行财产保全措施需要中止有关程序的，应当符合下列规定：

{\setlength{\parindent}{17.1bp}（1）人民法院应当将对专利申请权（或专利权）进行财产保全的民事裁定书及协助执行通知书送达专利局指定的接收部门，并提供人民法院的通讯地址、邮政编码和收件人姓名。\par}

{\setlength{\parindent}{17.1bp}（2）民事裁定书及协助执行通知书应当写明要求专利局协助执行的专利申请号（或专利号）、发明创造名称、申请人（或专利权人）的姓名或者名称、财产保全期限等内容。\par}

{\setlength{\parindent}{17.1bp}（3）要求协助执行财产保全的专利申请（或专利）处于有效期内。\par}

\guideparagraph{gpos.5.7.7.3.2.2}{7.3.2.2}{因协助执行财产保全而中止的审核及处理}{}{74.4bp}
专利局收到人民法院的民事裁定书和协助执行通知书后，应当按照本章第7.3.2.1节的规定进行审核，并按照下列情形处理：

{\setlength{\parindent}{17.1bp}（1）不符合规定的，应当向人民法院发出不予执行财产保全通知书，说明不执行中止的原因并继续原程序。\par}

{\setlength{\parindent}{17.1bp}（2）符合规定的，应当执行中止，并向人民法院和申请人（或专利权人）发出保全程序开始通知书，说明协助执行财产保全期限的起止日期（自收到民事裁定书之日起），并对专利权的财产保全予以公告。\par}

{\setlength{\parindent}{17.1bp}（3）对已执行财产保全的不得重复进行保全。执行中止后，其他人民法院又要求协助执行财产保全的，可以轮候保全。专利局应当进行轮候登记，对轮候登记在先的，自前一保全结束之日起轮候保全开始。\par}

对于处在无效宣告程序中的专利，专利局的流程管理部门还应当将执行中止的决定通知复审和无效审理部门，由复审和无效审理部门通知无效宣告程序中的当事人。

\guidesubsection{gpos.5.7.7.4}{7.4}{中止的期限}{}{102.0bp}
\guidesubsubsection{gpos.5.7.7.4.1}{7.4.1}{权属纠纷的当事人请求中止的期限}{\guidenote{细则 103.3}}{0bp}
对于专利申请权（或专利权）权属纠纷的当事人提出的中止请求，中止期限一般不得超过一年，即自中止请求之日起满一年的，该中止程序结束。

有关专利申请权（或专利权）权属纠纷在中止期限一年内未能结案，需要继续中止程序的，请求人应当在中止期满前请求延长中止期限，并提交权属纠纷受理部门出具的说明尚未结案原因的证明文件。中止程序可以延长一次，延长的期限不得超过六个月。不符合规定的，审查员应当发出延长期限审批通知书并说明不予延长的理由；符合规定的，审查员应当发出延长期限审批通知书，通知权属纠纷的双方当事人。

\guidesubsubsection{gpos.5.7.7.4.2}{7.4.2}{因协助执行财产保全而中止的期限}{\guidenote{细则 104}}{74.4bp}
对于人民法院要求专利局协助执行财产保全而执行中止程序的，按照民事裁定书及协助执行通知书写明的财产保全期限中止有关程序。

人民法院要求继续采取财产保全措施的，应当在中止期限届满前将继续保全的协助执行通知书送达专利局，经审核符合本章第7.3.2.1节规定的，中止期限予以续展。

\guidesubsubsection{gpos.5.7.7.4.3}{7.4.3}{涉及无效宣告程序的中止期限}{}{74.4bp}
对涉及无效宣告程序中的专利，应权属纠纷当事人请求的中止，中止期限不超过一年，中止期限届满专利局将自行恢复有关程序。

\guidesubsection{gpos.5.7.7.5}{7.5}{中止程序的结束}{}{102.0bp}
\guidesubsubsection{gpos.5.7.7.5.1}{7.5.1}{权属纠纷的当事人提出的中止程序的结束}{}{0bp}
\guidenote{细则 103.3}中止期限届满，专利局自行恢复有关程序，审查员应当向权属纠纷的双方当事人发出中止程序结束通知书，相关专利权已被宣告全部无效等情形除外。

对于尚在中止期限内的专利申请（或专利），地方知识产权管理部门作出的处理决定或者人民法院作出的判决产生法律效力之后（涉及权利人变更的，在办理著录项目变更手续之后），专利局应当结束中止程序。

专利局收到当事人、利害关系人、地方知识产权管理部门或者人民法院送交的调解书、裁定书或者判决书后，应当审查下列各项：

{\setlength{\parindent}{17.1bp}（1）文件是否有效，即是否是正式文本（正本或副本），是否是由有管辖权的机关作出的。\par}

{\setlength{\parindent}{17.1bp}（2）文件中记载的申请号（或专利号）、发明创造名称和权利人是否与请求结束中止程序的专利申请（或专利）中记载的内容一致。\par}

{\setlength{\parindent}{17.1bp}（3）文件是否已生效，如判决书的上诉期是否已满（调解书均没有上诉期）。当不能确定该文件是否已发生法律效力时，审查员应当给另一方当事人发出收到人民法院判决书的通知书，确认是否提起上诉；在指定的期限内未答复或者明确不上诉的，文件视为发生法律效力。提起上诉的，当事人应当提交上级人民法院出具的证明文件，原人民法院判决书不发生法律效力。\par}

文件不符合规定的，审查员应当向请求人发出视为未提出通知书，继续中止程序。文件符合规定并且未涉及权利人变更的，审查员应当发出中止程序结束通知书，通知双方当事人，恢复有关程序。

文件符合规定，但涉及权利人变更的，审查员应当发出办理手续补正通知书，通知取得权利一方的当事人在收到通知书之日起三个月内办理著录项目变更手续，并补办在中止程序中应办而未办的其他手续；取得权利一方的当事人办理有关手续后，审查员应当发出中止程序结束通知书，通知双方当事人，恢复有关程序。期满未办理有关手续的，视为放弃取得专利申请权（或专利权）的权利，审查员应当向取得权利的一方当事人发出视为放弃取得专利申请权或专利权的权利通知书，期满未办理恢复手续的，中止程序结束，审查员应当发出中止程序结束通知书，通知权属纠纷的双方当事人，恢复有关程序。

\guidesubsubsection{gpos.5.7.7.5.2}{7.5.2}{因人民法院要求协助执行财产保全的中止程序的结束}{\guidenote{细则 104}}{74.4bp}
中止期限届满，人民法院没有要求继续采取财产保全措施的，审查员应当发出中止程序结束通知书，通知人民法院和申请人（或专利权人），恢复有关程序，并对专利权保全解除予以公告。有轮候保全登记的，对轮候登记在先的，自前一保全结束之日起轮候保全开始，中止期限为民事裁定书及协助执行通知书写明的财产保全期限。审查员应当向前一个人民法院和申请人（或专利权人）发出中止程序结束通知书，向轮候登记在先的人民法院和申请人（或专利权人）发出保全程序开始通知书，说明协助执行财产保全期的起止日期，并对专利权的财产保全予以公告。

要求协助执行财产保全的人民法院送达解除保全通知书后，经审核符合规定的，审查员应当发出中止程序结束通知书，通知人民法院和申请人（或专利权人），恢复有关程序，并对专利权的保全解除予以公告。

\guidesection{gpos.5.7.8}{8.}{审查的顺序}{}{102.0bp}
\guidesubsection{gpos.5.7.8.1}{8.1}{一般原则}{}{0bp}
对于发明、实用新型和外观设计专利申请，一般应当按照申请提交的先后顺序启动初步审查；对于发明专利申请，在符合启动实质审查程序的其他条件前提下，一般应当按照提交实质审查请求书并缴纳实质审查费的先后顺序启动实质审查；另有规定的除外。

如果需要，对于技术内容、申请人或发明人相互关联的专利申请，可以合并审查。

\guidesubsection{gpos.5.7.8.2}{8.2}{优先审查}{}{74.4bp}
对涉及国家、地方政府重点发展或鼓励的产业，对国家利益或者公共利益具有重大意义的申请，或者在市场活动中具有一定需求的申请等，由申请人提出请求，经批准后，可以优先审查，并在随后的审查过程中予以优先处理。按照规定由其他相关主体提出优先审查请求的，依照规定处理。适用优先审查的具体情形由《专利优先审查管理办法》规定。

但是，同一申请人同日（仅指申请日）对同样的发明创造既申请实用新型又申请发明的，对于其中的发明专利申请一般不予优先审查。

\guidesubsection{gpos.5.7.8.3}{8.3}{延迟审查}{}{74.4bp}
申请人可以对专利申请提出延迟审查请求。

\guidenote{细则 56.2}发明专利延迟审查请求，应当由申请人在提出实质审查请求的同时提出，但发明专利申请延迟审查请求自实质审查请求生效之日起生效。延迟期限为自延迟审查请求生效之日起1 年、2年或者3年。

实用新型专利延迟审查请求，应当由申请人在提交实用新型专利申请的同时提出。延迟期限为自延迟审查请求生效之日起1年。

外观设计专利延迟审查请求，应当由申请人在提交外观设计申请的同时提出。延迟期限以月为单位，最长延迟期限为自延迟审查请求生效之日起36个月。

延迟期限届满后，专利申请将按顺序待审。必要时，专利局可以自行启动审查程序并通知申请人，申请人请求的延迟审查期限终止。

延迟期限届满前，申请人可以请求撤回延迟审查请求，符合规定的，延迟期限终止，专利申请将按顺序待审。

\guidesubsection{gpos.5.7.8.4}{8.4}{专利局自行启动}{}{74.4bp}
对于专利局自行启动实质审查的专利申请，可以优先处理。

\guidechapter{gpos.5.8}{第八章}{专利公报和单行本的编辑}{}{0bp}
\guidesection{gpos.5.8.1}{1.}{专利公报}{\guidenote{细则107}}{102.0bp}
\guidesubsection{gpos.5.8.1.1}{1.1}{专利公报的种类}{}{0bp}
专利局编辑出版的专利公报有发明专利公报、实用新型专利公报和外观设计专利公报。专利公报以电子公报形式、期刊形式或者以专利局规定的其他形式公布。电子公报在国家知识产权局政府网站上公布。三种专利公报按照年度计划定期分别出版。

\guidesubsection{gpos.5.8.1.2}{1.2}{专利公报的内容}{}{102.0bp}
\guidesubsubsection{gpos.5.8.1.2.1}{1.2.1}{发明专利公报}{}{0bp}
发明专利公报包括发明专利申请公布、国际专利申请公布、发明专利权授予、保密发明专利、发明专利事务、索引（申请公布索引、授权公告索引）。

\guideparagraph{gpos.5.8.1.2.1.1}{1.2.1.1}{发明专利申请公布}{}{74.4bp}
发明专利申请经初步审查合格后，自申请日（有优先权的，为优先权日）起满十八个月，即行公布。专利法实施细则第四十一条规定的作好公布专利申请文件印刷准备工作的时间一般为专利法第三十四条规定的十八个月前一个月。

发明专利申请人在初步审查合格前，要求提前公布其专利申请的，自初步审查合格之日起进行公布准备；在初步审查合格后，要求提前公布其专利申请的，自提前公布请求合格之日起进行公布准备。

在初步审查程序中被驳回、被视为撤回以及在作好公布准备之前申请人主动撤回或确定保密的发明专利申请不予公布。

发明专利申请公布的内容包括：著录事项、摘要和摘要附图，但说明书没有附图的，可以没有摘要附图。著录事项主要包括：国际专利分类号、申请号、公布号（出版号）、公布日、申请日、优先权事项、申请人事项、发明人事项、专利代理事项、发明名称等。

\guideparagraph{gpos.5.8.1.2.1.2}{1.2.1.2}{发明专利权授予}{}{74.4bp}
发明专利申请人根据专利局作出的授予专利权通知和办理登记手续通知，按时缴纳授予专利权当年的年费和其他有关费用后，该专利申请进入授权公告准备，并予以公告。

发明专利权授予公告的内容包括：国际专利分类号、专利号、授权公告号（出版号）、申请日、授权公告日、优先权事项、专利权人事项、发明人事项、专利代理事项、发明名称等著录事项。

\guideparagraph{gpos.5.8.1.2.1.3}{1.2.1.3}{保密发明专利和国防发明专利}{}{74.4bp}
保密发明专利只公告保密专利权的授予和保密专利的解密，保密专利公告的著录事项包括：专利号、申请日、授权公告日等。

保密发明专利解密后，在专利公报的解密栏中予以公告，出版单行本。

国防发明专利权的授予和解密的公告参照上述规定执行。

\guideparagraph{gpos.5.8.1.2.1.4}{1.2.1.4}{发明专利事务}{}{74.4bp}
发明专利事务公布专利局对发明专利申请和发明专利作出的决定和通知。包括：实质审查请求的生效，专利局对专利申请自行进行实质审查的决定，发明专利申请公布后的驳回，发明专利申请公布后的撤回，发明专利申请公布后的视为撤回，视为放弃取得专利权，专利权的全部（或部分）无效宣告，专利权的终止，专利权的主动放弃，专利申请（或专利）权利的恢复，专利权期限的补偿，专利申请权、专利权的转移，专利实施的强制许可，专利实施许可合同的备案，专利权的质押、保全及其解除，专利实施的开放许可，专利权人的姓名或者名称、地址等著录事项的变更，文件的公告送达，专利局的更正，其他有关事项等。

\guideparagraph{gpos.5.8.1.2.1.5}{1.2.1.5}{索引}{}{74.4bp}
发明索引分申请公布索引和授权公告索引两种。每种索引又分国际分类号索引、申请号索引（或者专利号索引）、申请人索引（或者专利权人索引）和公布号/申请号（授权公告号/专利号）对照表索引。

\guidesubsubsection{gpos.5.8.1.2.2}{1.2.2}{实用新型专利公报}{}{74.4bp}
实用新型专利公报包括实用新型专利权授予、保密实用新型专利、实用新型专利事务和授权公告索引。

\guideparagraph{gpos.5.8.1.2.2.1}{1.2.2.1}{实用新型专利权授予}{}{74.4bp}
实用新型专利申请人根据专利局作出的授予专利权通知和办理登记手续通知，按时缴纳授予专利权当年的年费和其他有关费用后，该专利申请进入授权公告准备，并予以公告。

实用新型专利权授予公告的内容包括：著录事项、摘要和摘要附图。著录事项主要包括：国际专利分类号、专利号、授权公告号（出版号）、申请日、授权公告日、优先权事项、专利权人事项、发明人事项、专利代理事项、实用新型名称。

申请人在申请时对同样的发明创造已申请发明专利作出说明的，应予以公告。

\guideparagraph{gpos.5.8.1.2.2.2}{1.2.2.2}{保密实用新型专利}{}{74.4bp}
保密实用新型专利只公告保密专利权的授予和保密专利的解密，保密专利公告的著录事项包括：专利号、申请日、授权公告日等。

保密实用新型专利解密后，在专利公报的解密栏中予以公告，出版单行本。

\guideparagraph{gpos.5.8.1.2.2.3}{1.2.2.3}{实用新型专利事务}{}{74.4bp}
实用新型专利事务公布专利局对实用新型专利申请和实用新型专利作出的决定和通知。包括：专利权的全部（或部分）无效宣告，专利权的终止，专利权的主动放弃，避免重复授权放弃实用新型专利权，专利权的恢复，专利权的转移，专利实施的强制许可，专利实施许可合同的备案，专利权的质押、保全及解除，专利实施的开放许可，专利权人的姓名或者名称、地址等著录事项的变更，文件的公告送达，专利局的更正，其他有关事项等。

\guideparagraph{gpos.5.8.1.2.2.4}{1.2.2.4}{授权公告索引}{}{74.4bp}
实用新型授权公告索引包括国际专利分类号索引、专利号索引、专利权人索引和授权公告号/专利号对照表索引。

\guidesubsubsection{gpos.5.8.1.2.3}{1.2.3}{外观设计专利公报}{}{74.4bp}
外观设计专利公报包括外观设计专利权的授予、外观设计专利事务和授权公告索引。

\guideparagraph{gpos.5.8.1.2.3.1}{1.2.3.1}{外观设计专利权授予}{}{74.4bp}
外观设计专利申请人根据专利局作出的授予专利权通知和办理登记手续通知，按时缴纳授予专利权当年的年费和其他有关费用后，该专利申请进入授权公告准备，并予以公告。

外观设计专利权授予公告的内容包括：著录事项、外观设计专利的一幅图片或者照片。著录事项主要包括：分类号、专利号、授权公告号（出版号）、申请日、授权公告日、优先权事项、专利权人事项、设计人事项、专利代理事项、使用该外观设计的产品名称等。

\guideparagraph{gpos.5.8.1.2.3.2}{1.2.3.2}{外观设计专利事务}{}{74.4bp}
外观设计专利事务公布专利局对外观设计专利申请和外观设计专利作出的决定和通知。包括：专利权的全部（或部分）无效宣告，专利权的终止，专利权的主动放弃，专利权的恢复，专利权的转移，专利实施许可合同的备案，专利权的质押、保全及其解除，专利实施的开放许可，专利权人的姓名或者名称、地址等著录事项的变更，文件的公告送达，专利局的更正，其他有关事项等。

\guideparagraph{gpos.5.8.1.2.3.3}{1.2.3.3}{授权公告索引}{}{74.4bp}
外观设计授权公告索引包括外观设计分类号索引、专利号索引、专利权人索引和授权公告号/专利号对照表索引。

\guidesubsection{gpos.5.8.1.3}{1.3}{专利公报的编辑}{}{102.0bp}
\guidesubsubsection{gpos.5.8.1.3.1}{1.3.1}{申请文件的编辑}{}{0bp}
用于公布的发明专利申请文件以及用于授权公告的发明专利申请文件、实用新型专利申请文件或外观设计专利申请文件应当符合制版要求，著录事项应当与公布准备或授权公告准备时专利申请文档记载的内容一致。

发明专利申请公布或发明专利权及实用新型专利权的授予按照国际专利分类号顺序编辑，主分类号相同的按照申请号顺序编辑。

外观设计专利权的授予按照外观设计分类号顺序编辑，分类号相同的按照申请号顺序编辑。

专利公报每一版面分左右两栏，自上而下，自左至右连续编排。

\guidesubsubsection{gpos.5.8.1.3.2}{1.3.2}{事务部分的编辑}{}{74.4bp}
各种专利公报事务部分编辑的原则：

{\setlength{\parindent}{17.1bp}（1）授予专利权公告之前专利局对实用新型和外观设计专利申请的权利丧失作出的决定不予刊登；公布之前专利局对发明专利申请的权利丧失作出的决定不予刊登。\par}

{\setlength{\parindent}{17.1bp}（2）刊登专利局作出的各种已经生效的按照规定应当公告的决定。\par}

{\setlength{\parindent}{17.1bp}（3）同一期公报中公布两项以上相同事务时，按照主分类号顺序编辑，主分类号相同的按照申请号顺序编辑。\par}

\guideparagraph{gpos.5.8.1.3.2.1}{1.3.2.1}{实质审查请求的生效、专利局对发明专利申请自行进行实质审查的决定}{}{74.4bp}
本事务仅适用于发明专利申请。公布的项目包括：主分类号、专利申请号、申请日。

\guideparagraph{gpos.5.8.1.3.2.2}{1.3.2.2}{发明专利申请公布后的驳回、撤回和视为撤回}{}{74.4bp}
本事务仅适用于已公布的发明专利申请。公布的项目包括：主分类号、专利申请号、公布日。

\guideparagraph{gpos.5.8.1.3.2.3}{1.3.2.3}{发明专利申请视为放弃取得专利权}{}{74.4bp}
公布的项目包括：主分类号、专利申请号。

\guideparagraph{gpos.5.8.1.3.2.4}{1.3.2.4}{专利实施的强制许可}{}{74.4bp}
公布的项目包括：主分类号、专利号、授权公告日。

\guideparagraph{gpos.5.8.1.3.2.5}{1.3.2.5}{专利权的终止}{}{74.4bp}
公布的项目包括：主分类号、专利号、申请日、授权公告日。

\guideparagraph{gpos.5.8.1.3.2.6}{1.3.2.6}{专利权期限的补偿}{}{74.4bp}
专利权期限补偿公布的项目包括：主分类号、专利号、申请日、授权公告日、原专利权期限届满日、现专利权期限届满日。

药品专利权期限补偿公布的项目包括：主分类号、专利号、申请日、授权公告日、药品名称及经批准的适应症、原专利权期限届满日、现专利权期限届满日。

\guideparagraph{gpos.5.8.1.3.2.7}{1.3.2.7}{专利实施许可合同备案的生效、变更及注销}{}{74.4bp}
{\setlength{\parindent}{22.7bp}专利实施许可合同备案生效公布的项目包括：主分类号、专利号、备案号、让与人、受让人、发明名称、申请日、发明公布日、授权公告日、许可种类（独占、排他、普通）、备案日。\par}

专利实施许可合同备案变更公布的项目包括：主分类号、专利号、备案号、变更日、变更项（许可种类、让与人、受让人）及变更前后内容。

专利实施许可合同备案注销公布的项目包括：主分类号、专利号、备案号、让与人、受让人、许可合同备案解除日。

\guideparagraph{gpos.5.8.1.3.2.8}{1.3.2.8}{专利权质押登记、变更及注销}{}{74.4bp}
专利权质押登记公布的项目包括：主分类号、专利号、登记号、质押登记日、出质人、质权人、发明名称、申请日、授权公告日。

专利权质押登记变更公布的项目包括：主分类号、专利号、登记号、变更日、变更项（出质人、质权人）及变更前后内容。

专利权质押登记注销公布的项目包括：主分类号、专利号、登记号、出质人、质权人、申请日、授权公告日、注销日。

\guideparagraph{gpos.5.8.1.3.2.9}{1.3.2.9}{专利权的保全及其解除}{}{74.4bp}
保全公布的项目包括：主分类号、专利号、申请日、授权公告日、保全登记生效日。

保全解除公布的项目包括：主分类号、专利号、申请日、授权公告日、保全解除日。

\guideparagraph{gpos.5.8.1.3.2.10}{1.3.2.10}{专利申请权、专利权的转移}{}{74.4bp}
公布的项目包括：主分类号、专利申请号（专利号）、变更项目、变更前权利人、变更后权利人、登记生效日。

\guideparagraph{gpos.5.8.1.3.2.11}{1.3.2.11}{专利权的全部或者部分无效宣告}{}{74.4bp}
专利权全部无效宣告公布的项目包括：主分类号、专利号、授权公告日、无效宣告决定号、无效宣告决定日。

专利权部分无效宣告公布的内容包括：主分类号、专利号、授权公告日、无效宣告决定号、无效宣告决定日、维持有效的权利要求。

\guideparagraph{gpos.5.8.1.3.2.12}{1.3.2.12}{专利权的主动放弃}{}{74.4bp}
公布的项目包括：主分类号、专利号、申请日、授权公告日、放弃生效日。

\guideparagraph{gpos.5.8.1.3.2.13}{1.3.2.13}{避免重复授权放弃实用新型专利权}{}{74.4bp}
公布的项目包括：主分类号、专利号、申请日、授权公告日、放弃生效日。

\guideparagraph{gpos.5.8.1.3.2.14}{1.3.2.14}{权利的恢复}{}{74.4bp}
公布的项目包括：主分类号、专利申请号（专利号）、原决定名称、原决定公告日。

\guideparagraph{gpos.5.8.1.3.2.15}{1.3.2.15}{文件的公告送达}{}{74.4bp}
由于文件送交地址不清，专利局无法通知当事人在规定或者指定的期限内答复或者办理手续的，应当在通知事项栏中公布。公布的项目包括：主分类号、申请号、收件人、文件名称。

\guideparagraph{gpos.5.8.1.3.2.16}{1.3.2.16}{其他有关事项}{}{74.4bp}
各事务栏内未规定的其他需要公告的内容，在本栏内公布。

\guideparagraph{gpos.5.8.1.3.2.17}{1.3.2.17}{更正}{}{74.4bp}
专利局对专利公报上出现的印刷及其他错误，一经发现，应当在更正栏中及时更正。各种不同类型错误的更正分别公布。公布的项目包括：主分类号、申请号（或专利号）、原公告所在卷号、更正项目、更正前内容、更正后内容。

\guidesubsubsection{gpos.5.8.1.3.3}{1.3.3}{索引的编辑}{}{102.0bp}
\guideparagraph{gpos.5.8.1.3.3.1}{1.3.3.1}{分类号索引}{}{0bp}
对于发明和实用新型，按照国际专利分类号编辑；对于外观设计按照外观设计分类号编辑。

分类号索引按照分类号为序，分类号相同的以公布号或者授权公告号为序。

分类号索引的项目包括：分类号，公布号或者授权公告号。

\guideparagraph{gpos.5.8.1.3.3.2}{1.3.3.2}{申请号或者专利号索引}{}{74.4bp}
申请号或者专利号索引以申请号或者专利号为序。

申请号或者专利号索引的项目包括：申请号或者专利号，公布号或者授权公告号。

\guideparagraph{gpos.5.8.1.3.3.3}{1.3.3.3}{申请人或者专利权人索引}{}{74.4bp}
{\setlength{\parindent}{22.2bp}申请人或者专利权人索引以申请人或者专利权人的姓名或者名称的拼音顺序为序。第一汉字相同的以第二汉字的拼音顺序为序，以此类推。外文名称排列在最前面，并以字母顺序为序。申请人或者专利权人相同的，以公布号或者授权公告号为序。\par}

申请人或者专利权人索引的项目包括：申请人或者专利权人，公布号或者授权公告号。

\guideparagraph{gpos.5.8.1.3.3.4}{1.3.3.4}{公布号/申请号（授权公告号/专利号）索引}{}{74.4bp}
公布号/申请号（授权公告号/专利号）对照表索引以公布号（授权公告号）为序。

公布号/申请号（授权公告号/专利号）对照表索引的项目包括：公布号（授权公告号），申请号（专利号）。

\guidesection{gpos.5.8.2}{2.}{专利申请及专利单行本}{\guidenote{细则108}}{74.4bp}
专利申请及专利单行本定期与相应的专利公报同日出版。

\guidesubsection{gpos.5.8.2.1}{2.1}{单行本的种类}{}{74.4bp}
单行本的种类包括：发明专利申请单行本、发明专利单行本、实用新型专利单行本及外观设计专利单行本。

\guidesubsection{gpos.5.8.2.2}{2.2}{单行本的内容}{}{102.0bp}
\guidesubsubsection{gpos.5.8.2.2.1}{2.2.1}{发明专利申请单行本}{}{0bp}
发明专利申请单行本的文献种类代码为“A”。包括：扉页、权利要求书、说明书（说明书有附图的，包含说明书附图）。

扉页由著录事项、摘要、摘要附图组成，说明书无附图的，则没有摘要附图。其内容应当与同一天出版的专利公报中相应专利申请的内容一致。

权利要求书、说明书及其附图，应当以审查员作出的发明专利申请初步审查合格通知书中指明的文本为准。

\guidesubsubsection{gpos.5.8.2.2.2}{2.2.2}{发明专利单行本}{}{74.4bp}
发明专利单行本的文献种类代码为“B”。包括：扉页、权利要求书、说明书（说明书有附图的，包含说明书附图）。

扉页由著录事项、摘要、摘要附图组成，说明书无附图的，则没有摘要附图。其内容比同日出版的专利公报中相应发明专利的内容增加审查员项和对比文件项。

权利要求书、说明书及其附图应当以审查员作出的授予专利权通知书中指明的文本为准。

发明专利权授予之后，在无效宣告程序中权利要求书需要修改后才能维持专利权的，应当再次出版该修改后的权利要求书，其文献种类代码依次为“C1—C7”，并标明修改后的权利要求书的公告日。

\guidesubsubsection{gpos.5.8.2.2.3}{2.2.3}{实用新型专利单行本}{}{74.4bp}
实用新型专利单行本的文献种类代码为“U”。包括：扉页、权利要求书、说明书和说明书附图。

扉页由著录事项、摘要和摘要附图组成，其内容应当与同日出版的实用新型专利公报中相应实用新型专利的内容一致。

权利要求书、说明书及其附图，应当以审查员作出的授予专利权通知书中指明的文本为准。

实用新型专利权授予之后，在无效宣告程序中权利要求书需要修改后才能维持专利权的，应当再次出版该修改后的权利要求书，其文献种类代码依次为“Y1—Y7”，并标明修改后的权利要求书的公告日。

\guidesubsubsection{gpos.5.8.2.2.4}{2.2.4}{外观设计专利单行本}{}{74.4bp}
外观设计专利单行本的文献种类代码为“S”。包括：扉页、彩色外观设计图片或者照片以及简要说明。

扉页由著录事项、一幅外观设计图片或者照片组成，其内容应当与同一天出版的外观设计专利公报中相应的外观设计专利内容一致。

彩色图片或者照片以及简要说明应当以审查员作出的授予专利权通知书中指明的图片或者照片以及简要说明为准。

外观设计专利权授予之后，在无效宣告程序中图片或者照片需要修改后才能维持专利权的，应当再次出版该修改后的图片或者照片，其文献种类代码依次为“S1—S7”，并标明修改后的图片或者照片的公告日。

\guidesubsection{gpos.5.8.2.3}{2.3}{更正}{}{74.4bp}
\guidenote{细则64 及107}\guidenote[15.6bp]{（17）}专利局对发明专利申请单行本、发明专利单行本、实用新型专利单行本及外观设计专利单行本的错误，一经发现，应当及时更正，重新出版更正的专利申请或专利单行本，并在其扉页上作出标记。

\guidechapter{gpos.5.9}{第九章}{专利权的授予和终止}{}{0bp}
\guidesection{gpos.5.9.1}{1.}{专利权的授予}{}{129.6bp}
\guidesubsection{gpos.5.9.1.1}{1.1}{专利权授予的程序}{}{0bp}
\guidesubsubsection{gpos.5.9.1.1.1}{1.1.1}{授予专利权通知}{}{0bp}
\guidenote{法 39 及 40}发明专利申请经实质审查、实用新型和外观设计专利申请经初步审查，没有发现驳回理由的，专利局应当作出授予专利权的决定，颁发专利证书，并同时在专利登记簿和专利公报上予以登记和公告。专利权自公告之日起生效。

{\setlength{\parindent}{22.2bp}在授予专利权之前，专利局应当发出授予专利权的通知书。\par}

\guidesubsubsection{gpos.5.9.1.1.2}{1.1.2}{办理登记手续通知}{\guidenote{细则 60.1}}{74.4bp}
专利局发出授予专利权通知书的同时，应当发出办理登记手续通知书，申请人应当在收到该通知之日起两个月内办理登记手续。

\guidesubsubsection{gpos.5.9.1.1.3}{1.1.3}{登记手续}{\guidenote{细则 114}}{74.4bp}
申请人在办理登记手续时，应当按照办理登记手续通知书的要求缴纳授权当年的年费。

\guidesubsubsection{gpos.5.9.1.1.4}{1.1.4}{颁发专利证书、登记和公告授予专利权}{\guidenote{法 39 及 40}}{74.4bp}
\guidenote{细则 60.1}申请人在规定期限之内办理登记手续的，专利局应当颁发专利证书，并同时予以登记和公告，专利权自公告之日起生效。

申请人办理登记手续后，专利局应当制作专利证书，进行专利权授予登记和公告授予专利权决定的准备。专利证书按照相关规定制作并送达专利权人。

\guidesubsubsection{gpos.5.9.1.1.5}{1.1.5}{视为放弃取得专利权的权利}{\guidenote{细则 60.2}}{74.4bp}
专利局发出授予专利权的通知书和办理登记手续通知书后，申请人在规定期限内未按照本章第1.1.3节规定办理登记手续的，应当发出视为放弃取得专利权通知书。该通知书应当在办理登记手续期满一个月后作出，并指明恢复权利的法律程序。自该通知书发出之日起四个月期满，未办理恢复手续的，或者专利局作出不予恢复权利决定的，将专利申请进行失效处理。对于发明专利申请，视为放弃取得专利权的，还应当在专利公报上予以公告。

\guidesubsection{gpos.5.9.1.2}{1.2}{专利证书}{}{102.0bp}
\guidesubsubsection{gpos.5.9.1.2.1}{1.2.1}{专利证书的构成}{}{0bp}
专利证书应当记载与专利权有关的重要著录事项、国家知识产权局印记、局长签字、授权公告日和授权公告号等。

著录事项包括：专利证书号（顺序号）、发明创造名称、专利号（即申请号）、专利申请日、发明人或者设计人姓名、专利权人姓名或者名称、第一专利权人地址、该专利申请日时的发明人或者设计人姓名和该专利申请日时的申请人姓名或者名称等。当一件专利的著录事项过长，在一页上记载有困难的，可以增加附页。

\guidesubsubsection{gpos.5.9.1.2.2}{1.2.2}{专利证书的更换}{}{74.4bp}
专利权权属纠纷经地方知识产权管理部门调解或者人民法院调解或者判决后，专利权归还请求人的，在该调解或者判决发生法律效力后，当事人可以在办理变更专利权人手续合格后，请求专利局更换专利证书。专利权终止后，专利局不再更换专利证书。因专利权的转移、专利权人更名发生专利权人姓名或者名称变更的，均不予更换专利证书。

专利局发出的原纸质专利证书破损损坏的，专利权人可以请求更换专利证书。专利局可以重新制作电子专利证书发送给当事人，更换后的证书应当与原专利证书的内容一致。

\guidesubsubsection{gpos.5.9.1.2.3}{1.2.3}{专利证书错误的更正}{}{74.4bp}
专利证书中存在错误时，专利权人可以请求专利局更正。专利局经核实存在错误的，将原专利证书公告作废，颁发更正后的专利证书。

\guidesubsection{gpos.5.9.1.3}{1.3}{专利登记簿}{}{102.0bp}
\guidesubsubsection{gpos.5.9.1.3.1}{1.3.1}{专利登记簿的格式}{}{0bp}
\guidenote{细则 106}专利局授予专利权时应当建立专利登记簿。专利登记簿登记的内容包括：专利权的授予，专利申请权、专利权的转移，国防专利、保密专利的解密，专利权的无效宣告，专利权的终止，专利权的恢复，专利权期限的补偿，专利权的质押、保全及其解除，专利实施许可合同的备案，专利实施的开放许可，专利实施的强制许可以及专利权人姓名或者名称、国籍、地址的变更。

上述事项一经作出即在专利登记簿中记载，专利登记簿登记的事项以数据形式储存于数据库中，制作专利登记簿副本时，按照规定的格式打印而成，加盖证件专用章后生效。

\guidesubsubsection{gpos.5.9.1.3.2}{1.3.2}{专利登记簿的效力}{}{74.4bp}
授予专利权时，专利登记簿与专利证书上记载的内容是一致的，在法律上具有同等效力；专利权授予之后，专利的法律状态的变更仅在专利登记簿上记载，由此导致专利登记簿与专利证书上记载的内容不一致的，以专利登记簿上记载的法律状态为准。

\guidesubsubsection{gpos.5.9.1.3.3}{1.3.3}{专利登记簿副本}{}{74.4bp}
\guidenote{细则 145.1}专利登记簿副本依据专利登记簿制作。专利权授予公告之后，任何人都可以向专利局请求出具专利登记簿副本。请求出具专利登记簿副本的，应当提交办理文件副本请求书并缴纳相关费用。

专利局收到有关请求和费用后，应当制作专利登记簿副本，经与专利申请文档核对无误后，加盖证件专用章后发送请求人。

\guidesection{gpos.5.9.2}{2.}{根据专利法第四十二条第二款的专利权期限补偿}{\guidenote{法 42.2}}{74.4bp}
根据专利法第四十二条第二款的规定，自发明专利申请日起满四年，且自实质审查请求之日起满三年后授予发明专利权的，专利局应专利权人的请求，就发明专利在授权过程中的不合理延迟给予专利权期限补偿，但由申请人引起的不合理延迟除外。

\guidenote{细则 78.4}同一申请人同日对同样的发明创造既申请实用新型专利又申请发明专利，依照专利法实施细则第四十七条第四款的规定取得发明专利权的，该发明专利权的期限不适用专利法第四十二条第二款的规定。

\guidesubsection{gpos.5.9.2.1}{2.1}{请求的提出}{}{74.4bp}
\guidenote{细则 77}专利权期限补偿请求应当由专利权人提出。专利权人请求给予专利权期限补偿的，应当自专利授权公告之日起三个月内向专利局提出请求，并且缴纳相应费用。

已委托专利代理机构的，专利权期限补偿请求应当由专利代理机构办理。专利权属于多个专利权人共有的，且未委托专利代理机构的，专利权期限补偿请求应当由代表人办理。

\guidesubsection{gpos.5.9.2.2}{2.2}{补偿期限的确定}{}{74.4bp}
\guidenote{细则 78.1 及.2}给予专利权期限补偿的，补偿期限按照发明专利在授权过程中不合理延迟的实际天数计算，该实际天数是指自发明专利申请日起满四年且自实质审查请求之日起满三年之日至公告授予专利权之日的间隔天数，减去合理延迟的天数和由申请人引起的不合理延迟的天数。

对于国际申请和分案申请，不合理延迟的实际天数是指自国际申请进入中国国家阶段的日期或分案申请递交日起满四年且自实质审查请求之日起满三年至公告授予专利权之日的间隔天数，减去合理延迟的天数和由申请人引起的不合理延迟的天数。

实质审查请求之日是指申请人依照专利法第三十五条第一款规定提出实质审查请求并依照专利法实施细则第一百一十三条规定足额缴纳发明专利申请实质审查费之日。发明专利申请的实质审查请求之日早于专利法第三十四条所称公布之日的，专利法第四十二条第二款所称自实质审查请求之日起满三年应当自该公布日起计算。

\guidesubsubsection{gpos.5.9.2.2.1}{2.2.1}{授权过程中的合理延迟}{}{74.4bp}
\guidenote{细则 78.3}以下情形引起的延迟属于授权过程中的合理延迟：依照专利法实施细则第六十六条规定修改专利申请文件的复审程序、依照专利法实施细则第一百零三条规定的中止程序、依照专利法实施细则第一百零四条规定的保全措施、其他合理情形如行政诉讼程序等。

\guidesubsubsection{gpos.5.9.2.2.2}{2.2.2}{申请人引起的不合理延迟}{}{74.4bp}
\guidenote{细则 79}以下由申请人引起的不合理延迟，延迟的天数为：

{\setlength{\parindent}{17.1bp}（1）未在指定期限内答复专利局发出的通知引起的延迟，延迟的天数为期限届满日起至实际提交答复之日止。\par}

{\setlength{\parindent}{17.1bp}（2）申请延迟审查的，延迟的天数为实际延迟审查的天数。\par}

{\setlength{\parindent}{17.1bp}（3）援引加入引起的延迟，延迟的天数为根据专利法实施细则第四十五条引起的延迟天数。\par}

{\setlength{\parindent}{17.1bp}（4）请求恢复权利引起的延迟，延迟的天数为从原期限届满日起至同意恢复的恢复权利请求审批通知书发文日止。能证明该延迟是由专利局造成的除外。\par}

{\setlength{\parindent}{17.1bp}（5）自优先权日起30个月内办理进入中国国家阶段手续的国际申请，申请人未要求提前处理引起的延迟，延迟的天数为进入中国国家阶段之日起至自优先权日起满30 个月之日止。\par}

\guidesubsection{gpos.5.9.2.3}{2.3}{期限补偿请求的审批}{}{74.4bp}
\guidenote{细则 84}经审查后认为专利权期限补偿请求不符合期限补偿条件的，专利局应当给予请求人至少一次陈述意见和/或补正文件的机会。对于此后仍然不符合期限补偿条件的，应当作出不予期限补偿的决定。

经审查后认为专利权期限补偿请求符合期限补偿条件的，专利局应当作出给予期限补偿的决定，告知期限补偿的天数。

\guidesubsection{gpos.5.9.2.4}{2.4}{登记和公告}{\guidenote[-11.5bp]{细则 84}}{74.4bp}
专利局作出给予专利权期限补偿的决定后，应当将有关事项在专利登记簿上登记并在专利公报上公告。

\guidesection{gpos.5.9.3}{3.}{根据专利法第四十二条第三款的专利权期限补偿}{}{74.4bp}
\guidenote{法 42.3}根据专利法第四十二条第三款和专利法实施细则第八十条至第八十四条的规定，对于国务院药品监督管理部门批准上市的创新药和符合规定的改良型新药，应专利权人的请求，专利局可以对符合条件的发明专利给予药品专利权期限补偿，以弥补在专利权有效期内该新药上市审评审批占用的时间。

\guidesubsection{gpos.5.9.3.1}{3.1}{补偿条件}{}{74.4bp}
\guidenote{细则 81}请求药品专利权期限补偿应当满足以下条件：

{\setlength{\parindent}{17.1bp}（1）请求补偿的专利授权公告日应当早于药品上市许可申请获得批准之日；\par}

{\setlength{\parindent}{17.1bp}（2）提出补偿请求时，该专利权处于有效状态；\par}

{\setlength{\parindent}{17.1bp}（3）该专利尚未获得过药品专利权期限补偿；\par}

{\setlength{\parindent}{17.1bp}（4）请求补偿专利的权利要求包括了获得上市许可的新药相关技术方案；\par}

{\setlength{\parindent}{17.1bp}（5）一个药品同时存在多项专利的，专利权人只能请求对其中一项专利给予药品专利权期限补偿；\par}

{\setlength{\parindent}{17.1bp}（6）一项专利同时涉及多个药品的，只能对一个药品就该专利提出药品专利权期限补偿请求。\par}

\guidesubsection{gpos.5.9.3.2}{3.2}{请求的提出}{}{74.4bp}
\guidenote{细则 81}药品专利权期限补偿请求应当由专利权人提出。专利权人与药品上市许可持有人不一致的，应当征得药品上市许可持有人书面同意。

专利权人请求药品专利权期限补偿的，应当自药品在中国获得上市许可之日起三个月内向专利局提出请求，并且缴纳相应费用。对于获得附条件上市许可的药品，应当自在中国获得正式上市许可之日起三个月内向专利局提出请求，但补偿期限的计算以获得附条件上市许可之日为准。

已委托专利代理机构的，药品专利权期限补偿请求应当由专利代理机构办理。专利权属于多个专利权人共有的，且未委托专利代理机构的，药品专利权期限补偿请求应当由代表人办理。

\guidesubsection{gpos.5.9.3.3}{3.3}{证明材料}{}{74.4bp}
提出药品专利权期限补偿请求时，请求人还应当提交如下材料：

{\setlength{\parindent}{17.1bp}（1）专利权人与药品上市许可持有人不一致的，应当提交药品上市许可持有人的书面同意书等材料；\par}

{\setlength{\parindent}{17.1bp}（2）用于确定药品专利权期限补偿期间专利保护范围的相关技术资料，例如请求对制备方法专利进行期限补偿的，应当提交国务院药品监督管理部门核准的药品生产工艺资料；\par}

{\setlength{\parindent}{17.1bp}（3）专利局要求的其他证明材料。\par}

请求人应当在请求中说明药品名称、药品注册分类、批准的适应症和请求给予期限补偿的专利号，指定与获得上市许可新药相关的权利要求，结合证明材料具体说明指定权利要求包括了新药相关技术方案的理由以及请求补偿期限的计算依据，并明确药品专利权期限补偿期间保护的技术方案。

\guidesubsection{gpos.5.9.3.4}{3.4}{适用范围}{\guidenote{法 42.3}}{74.4bp}
\guidenote{细则 80}根据专利法第四十二条第三款及专利法实施细则第八十条的规定，针对国务院药品监督管理部门批准上市的创新药和符合本章规定的改良型新药，对于其中药物活性物质的产品发明专利、制备方法发明专利或者医药用途发明专利，可以给予药品专利权期限补偿。创新药和改良型新药的含义依照有关法律法规并按照国务院药品监督管理部门的相关规定确定。

可以给予期限补偿的改良型新药限于国务院药品监督管理部门颁发的药品注册证书中记载为以下类别的改良型新药：

{\setlength{\parindent}{17.1bp}（1）化学药品第2.1类中对已知活性成份成酯，或者对已知活性成份成盐的药品；\par}

{\setlength{\parindent}{17.1bp}（2）化学药品第2.4类，即含有已知活性成份的新适应症的药品；\par}

{\setlength{\parindent}{17.1bp}（3）预防用生物制品第2.2类中对疫苗菌毒种改进的疫苗；\par}

{\setlength{\parindent}{17.1bp}（4）治疗用生物制品第2.2类中增加新适应症的生物制品；\par}

{\setlength{\parindent}{17.1bp}（5）中药第2.3类，即增加功能主治的中药。\par}

\guidesubsection{gpos.5.9.3.5}{3.5}{指定权利要求是否包括新药相关技术方案的审查}{}{74.4bp}
新药相关技术方案应当以国务院药品监督管理部门批准的新药的结构、组成及其含量，批准的生产工艺和适应症为准。指定权利要求未包括获得上市许可的新药相关技术方案的，不予期限补偿。

\guidenote{细则 83}药品专利权期限补偿期间内，该专利的保护范围限于国务院药品监督管理部门批准上市的新药，且限于该新药经批准的适应症相关技术方案；在保护范围内，专利权人享有的权利和承担的义务与专利权期限补偿前相同。产品权利要求的保护范围仅限于用于经批准的适应症的上市新药产品，医药用途权利要求的保护范围仅限于上市新药产品的经批准的适应症，制备方法权利要求的保护范围仅限于用于经批准的适应症的上市新药产品在国务院药品监督管理部门备案的生产工艺。

\guidesubsection{gpos.5.9.3.6}{3.6}{补偿期限的确定}{}{74.4bp}
\guidenote{细则 82}给予药品专利权期限补偿的，补偿期限按照该专利申请日至该新药在中国获得上市许可之日的间隔天数减去5年。该补偿期限不超过5年，且该药品上市许可申请批准后总有效专利权期限不超过14年。

\guidesubsection{gpos.5.9.3.7}{3.7}{期限补偿请求的审批}{}{74.4bp}
\guidenote{细则 84}经审查后认为药品专利权期限补偿请求不符合期限补偿条件的，专利局应当给予请求人至少一次陈述意见和/或补正文件的机会。对于此后仍然不符合期限补偿条件的，应当作出不予期限补偿的决定。

经审查认为应当给予药品专利权期限补偿的，如果专利权人已经提出专利权期限补偿请求但专利局尚未作出审批决定，审查员应当等待专利权期限补偿请求的审批决定作出以后，再确定给予药品专利权期限补偿的时间；如果专利权人尚未提出专利权期限补偿请求，且其自专利授权公告之日起三个月期限尚未届满，审查员应当等待专利权期限补偿请求的时限届满以后，再确定给予药品专利权期限补偿的时间，但专利权人明确表示放弃提出专利权期限补偿请求的除外。

经审查后认为药品专利权期限补偿请求符合期限补偿条件的，专利局应当作出给予期限补偿的决定，告知期限补偿的天数。

\guidesubsection{gpos.5.9.3.8}{3.8}{登记和公告}{\guidenote{细则 84}}{74.4bp}
专利局作出给予药品专利权期限补偿的决定后，应当将有关事项在专利登记簿上登记并在专利公报上公告。

\guidesection{gpos.5.9.4}{4.}{专利权的终止}{}{102.0bp}
\guidesubsection{gpos.5.9.4.1}{4.1}{专利权期满终止}{}{0bp}
\guidenote{法 42}发明专利权的期限为二十年，实用新型专利权期限为十年，外观设计专利权期限为十五年，均自申请日起计算。例如，一件实用新型专利的申请日是1999年9月6日，该专利的期限为1999年9月6日至2009年9月5日，专利权期满终止日为2009年9月6日（遇节假日不顺延）。

发明专利权如存在专利法第四十二条第二款、第三款的期限补偿的，专利权期满终止日为期限补偿后的专利权期满终止日。例如，一件发明专利的申请日为2021年9月6日，该专利的期限为2021年9月6日至2041年9月5日。如其专利权期限补偿后的专利权期限届满日为2041年12月1日，则该发明专利的专利权期满终止日为2041年12月2日（遇节假日不顺延）。

专利权期满时应当及时在专利登记簿和专利公报上分别予以登记和公告，并进行失效处理。

\guidesubsection{gpos.5.9.4.2}{4.2}{专利权人没有按照规定缴纳年费的终止}{}{102.0bp}
\guidesubsubsection{gpos.5.9.4.2.1}{4.2.1}{年费}{\guidenote{细则 115}}{0bp}
授予专利权当年的年费应当在办理登记手续的同时缴纳，以后的年费应当在上一年度期满前缴纳。缴费期限届满日是申请日在该年的相应日。

\guideparagraph{gpos.5.9.4.2.1.1}{4.2.1.1}{年度}{}{74.4bp}
专利年度从申请日起算，与优先权日、授权日无关，与自然年度也没有必然联系。例如，一件专利申请的申请日是1999年6 月1 日，该专利申请的第一年度是1999 年6 月1 日至2000年5月31日，第二年度是2000年6月1日至2001年5月31日，以此类推。

\guideparagraph{gpos.5.9.4.2.1.2}{4.2.1.2}{应缴年费数额}{}{74.4bp}
各年度年费按照收费表中规定的数额缴纳。例如，一件专利申请的申请日是1997年6月3日，如果该专利申请于2001年8月1日被授予专利权（授予专利权公告之日），申请人在办理登记手续时已缴纳了第五年度年费，那么该专利权人最迟应当在2002年6月3日按照第六年度年费标准缴纳第六年度年费。

\guideparagraph{gpos.5.9.4.2.1.3}{4.2.1.3}{滞纳金}{}{74.4bp}
专利权人未按时缴纳年费（不包括授予专利权当年的年费）或者缴纳的数额不足的，可以在年费期满之日起六个月内补缴，补缴时间超过规定期限但不足一个月时，不缴纳滞纳金。补缴时间超过规定时间一个月或以上的，缴纳按照下述计算方法算出的相应数额的滞纳金：

{\setlength{\parindent}{17.1bp}（1）超过规定期限一个月（不含一整月）至两个月（含两个整月）的，缴纳数额为全额年费的5\%。\par}

{\setlength{\parindent}{17.1bp}（2）超过规定期限两个月至三个月（含三个整月）的，缴纳数额为全额年费的10\%。\par}

{\setlength{\parindent}{17.1bp}（3）超过规定期限三个月至四个月（含四个整月）的，缴纳数额为全额年费的15\%。\par}

{\setlength{\parindent}{17.1bp}（4）超过规定期限四个月至五个月（含五个整月）的，缴纳数额为全额年费的20\%。\par}

{\setlength{\parindent}{17.1bp}（5）超过规定期限五个月至六个月的，缴纳数额为全额年费的25\%。\par}

凡在六个月的滞纳期内补缴年费或者滞纳金不足需要再次补缴的，应当依照再次补缴年费或者滞纳金时所在滞纳金时段内的滞纳金标准，补足应当缴纳的全部年费和滞纳金。例如，年费滞纳金5\%的缴纳时段为5月10日至6月10日，滞纳金为45元，但缴费人仅缴了25元。缴费人在6月15日补缴滞纳金时，应当依照再次缴费日所对应的滞纳期时段的标准10\%缴纳。该时段滞纳金金额为90元，还应当补缴65元。

凡因年费和/或滞纳金缴纳逾期或者不足而造成专利权终止的，在恢复程序中，除补缴年费之外，还应当缴纳或者补足全额年费25\%的滞纳金。

\guidesubsubsection{gpos.5.9.4.2.2}{4.2.2}{终止}{}{74.4bp}
专利年费滞纳期满仍未缴纳或者缴足专利年费或者滞纳金的，审查员应当发出专利权终止通知书。专利权人未启动恢复程序或者恢复权利请求未被批准的，专利局应当在终止通知书发出四个月后，进行失效处理，并在专利公报上公告。

专利权自应当缴纳年费期满之日起终止。

\guidesubsection{gpos.5.9.4.3}{4.3}{专利权人放弃专利权}{\guidenote{法 44.1（2）}}{74.4bp}
授予专利权后，专利权人随时可以主动要求放弃专利权，专利权人放弃专利权的，应当提交放弃专利权声明，并附具全体专利权人签字或者盖章同意放弃专利权的证明材料，或者仅提交由全体专利权人签字或者盖章的放弃专利权声明。委托专利代理机构的，放弃专利权的手续应当由专利代理机构办理，并附具全体申请人签字或者盖章的同意放弃专利权声明。主动放弃专利权的声明不得附有任何条件。放弃专利权只能放弃一件专利的全部，放弃部分专利权的声明视为未提出。

放弃专利权声明经审查，不符合规定的，审查员应当发出视为未提出通知书；符合规定的，审查员应当发出手续合格通知书，并将有关事项分别在专利登记簿和专利公报上登记和公告。放弃专利权声明的生效日为手续合格通知书的发文日，放弃的专利权自该日起终止。专利权人无正当理由不得要求撤销放弃专利权的声明。除非在专利权非真正拥有人恶意要求放弃专利权后，专利权真正拥有人（应当提供生效的法律文书来证明）可要求撤销放弃专利权声明。

\guidenote{细则 47.5}申请人依据专利法第九条第一款和专利法实施细则第四十七条第四款声明放弃实用新型专利权的，专利局在公告授予发明专利权时对放弃实用新型专利权的声明予以登记和公告。在无效宣告程序中声明放弃实用新型专利权的，专利局及时登记和公告该声明。放弃实用新型专利权声明的生效日为发明专利权的授权公告日，放弃的实用新型专利权自该日起终止。

\guidechapter{gpos.5.10}{第十章}{专利权评价报告}{}{0bp}
\guidesection{gpos.5.10.1}{1.}{引言}{}{74.4bp}
\guidenote{法 66.2}\guidenote[15.7bp]{细则 62.1}专利法第六十六条第二款规定，专利侵权纠纷涉及实用新型专利或者外观设计专利的，人民法院或者管理专利工作的部门可以要求专利权人或者利害关系人出具由国家知识产权局作出的专利权评价报告；专利权人、利害关系人或者被控侵权人也可以主动出具专利权评价报告。

国家知识产权局根据请求，对相关实用新型专利或者外观设计专利进行检索，并就该专利是否符合专利法及其实施细则规定的授权条件进行分析和评价，作出专利权评价报告。

专利权评价报告是人民法院或者管理专利工作的部门审理、处理专利侵权纠纷的证据，主要用于人民法院或者管理专利工作的部门确定是否需要中止相关程序。专利权评价报告不是行政决定，因此请求人不能就此提起行政复议和行政诉讼。

以实用新型或者外观设计专利权办理转让、出质登记和专利实施许可合同备案的,必要时,国家知识产权局可以要求提交专利权评价报告。

\guidesection{gpos.5.10.2}{2.}{专利权评价报告请求的形式审查}{}{74.4bp}
国家知识产权局收到专利权人、利害关系人或者被控侵权人提交的专利权评价报告请求后，应当进行形式审查。

\guidesubsection{gpos.5.10.2.1}{2.1}{专利权评价报告请求的主体及时机}{\guidenote{细则 62.1 及 63}}{74.4bp}
授予实用新型或者外观设计专利权的决定公告后，专利权人、利害关系人或者被控侵权人可以请求国家知识产权局作出专利权评价报告。申请人也可以在办理专利权登记手续时请求国家知识产权局作出专利权评价报告。

实用新型或者外观设计专利权属于多个专利权人共有的，请求人可以是部分专利权人。

利害关系人是指有权根据专利法第六十五条的规定就专利侵权纠纷向人民法院起诉或者请求管理专利工作的部门处理的人，例如专利实施独占许可合同的被许可人和由专利权人授予起诉权的专利实施普通许可合同的被许可人。

收到专利权人发出的律师函、电商平台投诉通知书等的单位或者个人属于被控侵权人，可以请求国家知识产权局作出专利权评价报告。

不符合上述规定的，专利权评价报告请求视为未提出。

\guidesubsection{gpos.5.10.2.2}{2.2}{专利权评价报告请求的客体}{\guidenote{细则 62.1}}{74.4bp}
专利权评价报告请求的客体应当是已经授权公告的实用新型专利或者外观设计专利，包括已经终止或者放弃的实用新型专利或者外观设计专利。针对下列情形提出的专利权评价报告请求视为未提出：

{\setlength{\parindent}{17.1bp}（1）未授权公告的实用新型专利申请或者外观设计专利申请，申请人在办理登记手续时提交专利权评价报告请求的除外；\par}

{\setlength{\parindent}{17.1bp}（2）已被宣告全部无效的实用新型专利或者外观设计专利；\par}

{\setlength{\parindent}{17.1bp}（3）已作出专利权评价报告的实用新型专利或者外观设计专利。\par}

\guidesubsection{gpos.5.10.2.3}{2.3}{专利权评价报告请求书}{\guidenote{细则 62.2 及.3}}{74.4bp}
在请求作出专利权评价报告时，请求人应当提交专利权评价报告请求书及相关的文件。

{\setlength{\parindent}{17.1bp}（1）专利权评价报告请求书应当采用国家知识产权局规定的表格。请求书中应当写明实用新型专利或者外观设计专利的申请号或者专利号、发明创造名称、申请人或者专利权人、请求人名称或者姓名。每一请求应当限于一件实用新型或者外观设计专利。\par}

{\setlength{\parindent}{17.1bp}（2）请求人是利害关系人的，在提出专利权评价报告请求的同时应当提交相关证明文件。例如，请求人是专利实施独占许可合同的被许可人的，应当提交与专利权人订立的专利实施独占许可合同或其复印件；请求人是专利权人授予起诉权的专利实施普通许可合同的被许可人的，应当提交与专利权人订立的专利实施普通许可合同或其复印件，以及专利权人授予起诉权的证明文件。如果所述专利实施许可合同已在国家知识产权局备案，请求人可以不提交专利实施许可合同，但应在请求书中注明。\par}

{\setlength{\parindent}{17.1bp}（3）请求人是被控侵权人的，在提交专利权评价报告请求的同时应当提交相关证明文件。例如，人民法院出具的立案类通知书或其复印件，专利行政执法部门出具的立案类通知书或其复印件，调解仲裁机构出具的立案类文件或其复印件，专利权人发出的律师函或其复印件，电商平台投诉通知书或其复印件等。\par}

专利权评价报告请求书不符合上述规定的，国家知识产权局应当通知请求人在指定期限内补正。

\guidesubsection{gpos.5.10.2.4}{2.4}{费用}{\guidenote{细则 110 、 111}}{74.4bp}
\guidenote{及 116.3}请求人自提出专利权评价报告请求之日起一个月内未缴纳或者未缴足专利权评价报告请求费的，专利权评价报告请求视为未提出。

\guidesubsection{gpos.5.10.2.5}{2.5}{委托手续}{}{74.4bp}
专利权评价报告请求的相关事务可以由请求人或者其委托的专利代理机构办理。对于根据专利法第十八条第一款规定应当委托专利代理机构的请求人，未按规定委托的，国家知识产权局应当通知请求人在指定期限内补正。

请求人委托专利代理机构办理专利权评价报告有关手续的，应当提交委托书，并在委托书中写明其委托权限仅限于办理专利权评价报告相关事务；委托手续不符合规定的，国家知识产权局应当要求请求人在指定期限内补正；期满未补正或者在指定期限内补正不符合规定的，视为未提出专利权评价报告请求。申请人或专利权人委托原专利代理机构办理专利权评价报告手续的，不需要提交委托书。

\guidesubsection{gpos.5.10.2.6}{2.6}{形式审查后的处理}{}{74.4bp}
{\setlength{\parindent}{17.1bp}\guidenote{细则 62.3}（1）专利权评价报告请求经形式审查不符合规定需要补正的，国家知识产权局应当发出补正通知书，要求请求人在收到通知书之日起十五日内补正；期满未补正或者在指定期限内补正但经两次补正后仍存在同样缺陷的，其请求视为未提出。\par}

{\setlength{\parindent}{17.1bp}（2）专利权评价报告请求视为未提出的，国家知识产权局应当发出视为未提出通知书，通知请求人。\par}

{\setlength{\parindent}{17.1bp}（3）专利权评价报告请求经形式审查合格的，应当及时转送给指定的作出专利权评价报告的部门。\par}

根据专利法实施细则第六十三条的规定，作出专利权评价报告前，多个请求人分别请求对同一件实用新型专利或者外观设计专利作出专利权评价报告的，国家知识产权局均予以受理，但仅作出一份专利权评价报告。

\guidesection{gpos.5.10.3}{3.}{专利权评价}{}{74.4bp}
作出专利权评价报告的部门在收到专利权评价报告请求书后，应当指派审查员按照本章的规定对该专利进行检索、分析和评价，作出专利权评价报告。

\guidesubsection{gpos.5.10.3.1}{3.1}{核查专利权评价报告请求书}{}{74.4bp}
审查员首先应当核查专利权评价报告请求书及其相关文件。发现不符合规定的，返回相应的部门处理，并说明理由。

\guidesubsection{gpos.5.10.3.2}{3.2}{专利权评价的内容}{}{102.0bp}
\guidesubsubsection{gpos.5.10.3.2.1}{3.2.1}{实用新型专利}{}{0bp}
实用新型专利权评价所涉及的内容包括：

{\setlength{\parindent}{17.1bp}（1）实用新型是否属于专利法第五条或者第二十五条规定的不授予专利权的情形，其评价标准适用本指南第二部分第一章的规定。\par}

{\setlength{\parindent}{17.1bp}（2）实用新型是否属于专利法第二条第三款规定的客体，其评价标准适用本指南第一部分第二章第6节的规定。\par}

{\setlength{\parindent}{17.1bp}（3）实用新型是否具备专利法第二十二条第四款规定的实用性，其评价标准适用本指南第二部分第五章第3节的规定。\par}

{\setlength{\parindent}{17.1bp}（4）实用新型专利的说明书是否按照专利法第二十六条第三款的要求充分公开了专利保护的主题，其评价标准适用本指南第二部分第二章第2.1节的规定。\par}

{\setlength{\parindent}{17.1bp}（5）实用新型是否具备专利法第二十二条第二款规定的新颖性，其评价标准适用本指南第四部分第六章第3节的规定。\par}

{\setlength{\parindent}{17.1bp}（6）实用新型是否具备专利法第二十二条第三款规定的创造性，其评价标准适用本指南第四部分第六章第4节的规定。\par}

{\setlength{\parindent}{17.1bp}（7）实用新型是否符合专利法第二十六条第四款的规定，其评价标准适用本指南第二部分第二章第3.2节的规定。\par}

{\setlength{\parindent}{17.1bp}（8）实用新型是否符合专利法实施细则第二十三条第二款的规定，其评价标准适用本指南第二部分第二章第3.1.2节的规定。\par}

{\setlength{\parindent}{17.1bp}（9）实用新型专利文件的修改是否符合专利法第三十三条的规定，其评价标准适用本指南第一部分第二章第8节和第二部分第八章第5.2节的规定。\par}

{\setlength{\parindent}{17.1bp}（10）分案的实用新型专利是否符合专利法实施细则第四十九条第一款的规定，其评价标准适用本指南第二部分第六章第3.2节的规定。\par}

{\setlength{\parindent}{17.1bp}（11）实用新型是否符合专利法第九条的规定，其评价标准适用本指南第二部分第三章第6节的规定。\par}

{\setlength{\parindent}{17.1bp}（12）实用新型是否符合专利法实施细则第十一条的规定，其评价标准适用《规范申请专利行为的规定》。\par}

\guidesubsubsection{gpos.5.10.3.2.2}{3.2.2}{外观设计专利}{}{74.4bp}
外观设计专利权评价所涉及的内容包括：

{\setlength{\parindent}{17.1bp}（1）外观设计是否属于专利法第五条或者第二十五条规定的不授予专利权的情形，其评价标准适用本指南第一部分第三章第6.1和6.2节的规定。\par}

{\setlength{\parindent}{17.1bp}（2）外观设计是否属于专利法第二条第四款规定的客体，其评价标准适用本指南第一部分第三章第7节的规定。\par}

{\setlength{\parindent}{17.1bp}（3）外观设计是否符合专利法第二十三条第一款的规定，其评价标准适用本指南第四部分第五章第5节的规定。\par}

{\setlength{\parindent}{17.1bp}（4）外观设计是否符合专利法第二十三条第二款的规定，其评价标准适用本指南第四部分第五章第6节的规定。\par}

{\setlength{\parindent}{17.1bp}（5）外观设计专利的图片或者照片是否符合专利法第二十七条第二款的规定，其评价标准适用本指南第一部分第三章第4节的规定。\par}

{\setlength{\parindent}{17.1bp}（6）外观设计专利文件的修改是否符合专利法第三十三条的规定，其评价标准适用本指南第一部分第三章第10 节的规定。\par}

{\setlength{\parindent}{17.1bp}（7）分案的外观设计专利是否符合专利法实施细则第四十九条第一款的规定，其评价标准适用本指南第一部分第三章第9.4.2节的规定。\par}

{\setlength{\parindent}{17.1bp}（8）外观设计是否符合专利法第九条的规定，其评价标准适用本指南第四部分第五章第8节的规定。\par}

{\setlength{\parindent}{17.1bp}（9）外观设计是否符合专利法实施细则第十一条的规定，其评价标准适用《规范申请专利行为的规定》。\par}

\guidesubsection{gpos.5.10.3.3}{3.3}{检索}{}{74.4bp}
一般情况下，作出实用新型专利权评价报告或者外观设计专利权评价报告前，都应当进行检索。

\guidesubsubsection{gpos.5.10.3.3.1}{3.3.1}{实用新型专利}{}{74.4bp}
检索应当针对实用新型专利的所有权利要求进行，但实用新型专利保护的主题属于下列情形之一的，审查员对该主题不必进行检索：

{\setlength{\parindent}{17.1bp}（1）不符合专利法第二条第三款的规定；\par}

{\setlength{\parindent}{17.1bp}（2）属于专利法第五条或者第二十五条规定的不授予专利权的情形；\par}

{\setlength{\parindent}{17.1bp}（3）不具备实用性；\par}

{\setlength{\parindent}{17.1bp}（4）说明书和权利要求书未对该主题作出清楚、完整的说明，以致于所属技术领域的技术人员不能实现。\par}

检索的具体要求可以参照本指南第二部分第七章。

\guidesubsubsection{gpos.5.10.3.3.2}{3.3.2}{外观设计专利}{}{74.4bp}
检索应当针对外观设计专利的图片或照片表示的所有产品外观设计进行，并考虑简要说明的内容。但外观设计专利保护的产品外观设计属于下列情形之一的，审查员不必对该产品外观设计进行检索：

{\setlength{\parindent}{17.1bp}（1）不符合专利法第二条第四款的规定；\par}

{\setlength{\parindent}{17.1bp}（2）属于专利法第五条或者第二十五条规定的不授予专利权的情形；\par}

{\setlength{\parindent}{17.1bp}（3）图片或者照片未清楚地显示要求专利保护的产品的外观设计。\par}

审查员应当检索外观设计专利在中国提出申请之日以前公开的外观设计。为了确定是否存在抵触申请，审查员应当检索在该外观设计专利的申请日之前向专利局提交、并且在该外观设计专利的申请日后公告的外观设计专利。为了确定是否存在重复授权，审查员还应当检索在该外观设计专利的申请日向专利局提交的、并且已经公告的外观设计专利。

\guidesection{gpos.5.10.4}{4.}{专利权评价报告}{}{74.4bp}
\guidenote{细则 63.1}国家知识产权局应当自收到合格的专利权评价报告请求书和请求费后两个月内作出专利权评价报告；但申请人在办理专利权登记手续时请求作出专利权评价报告且形式审查合格的，国家知识产权局应当自授权公告之日起两个月内作出专利权评价报告。

未发现被评价专利存在不符合专利法及其实施细则规定的授予专利权条件的，审查员应当在专利权评价报告中给出明确结论。

对于被评价专利存在不符合专利法及其实施细则规定的授予专利权条件的，审查员应当在专利权评价报告中根据专利法及其实施细则具体阐述评价意见，并给出该专利不符合专利法及其实施细则规定的授予专利权条件的明确结论。

专利权评价报告使用国家知识产权局统一制定的标准表格，作出后由审查员与审核员署名，并加盖“国家知识产权局专利权评价报告专用章”。

\guidesubsection{gpos.5.10.4.1}{4.1}{专利权评价报告的内容}{}{74.4bp}
专利权评价报告包括反映对比文件与被评价专利相关程度的表格部分，以及该专利是否符合专利法及其实施细则规定的授予专利权的条件的说明部分。

\guidesubsubsection{gpos.5.10.4.1.1}{4.1.1}{表格部分}{}{74.4bp}
对于实用新型专利权评价报告，其表格部分的填写要求参见本指南第二部分第七章第12节的规定。

对于外观设计专利权评价报告，其表格部分应当清楚地记载检索的领域、数据库、由检索获得的对比文件以及对比文件与外观设计专利的相关程度等内容。通常，采用下列符号表示对比文件与外观设计专利的关系：

X：单独导致外观设计专利不符合专利法第二十三条第一款或第二款规定的文件；

Y：与报告中其他文件结合导致外观设计专利不符合专利法第二十三条第二款规定的文件；

A：背景文件，即反映外观设计的部分设计特征或者有关的文件；

P：中间文件，其公开日在外观设计专利的申请日与所要求的优先权日之间的文件，或者会导致需要核实外观设计专利优先权的文件；

E：与外观设计专利相同或者实质相同的抵触申请文件；

R：任何单位或个人在申请日向专利局提交的、属于同样的发明创造的外观设计专利文件。

上述类型的文件中，符号X、Y和A表示对比文件与外观设计专利在内容上的相关程度；符号R和E同时表示对比文件与外观设计专利在时间上的关系和在内容上的相关程度；符号P表示对比文件与外观设计专利在时间上的关系，其后应附带标明文件内容相关程度的符号X、Y、E或A，它属于在未核实优先权的情况下所作的标记。

\guidesubsubsection{gpos.5.10.4.1.2}{4.1.2}{说明部分}{}{74.4bp}
说明部分应当记载和反映专利权评价的结论。对于不符合专利法及其实施细则规定的授予专利权条件的被评价专利，还应当给出明确、具体的评价意见。

{\setlength{\parindent}{17.1bp}（1）对于不符合专利法及其实施细则规定的授予专利权条件的实用新型专利，应当给出具体的评价说明，并明确结论，必要时应当引证对比文件。例如，对于不具备新颖性和/或创造性的权利要求，审查员应当逐一进行评述；对于多项从属权利要求，应当对其引用不同的权利要求时的技术方案分别进行评述；对于具有并列选择方案的权利要求，应当对各选择方案分别进行评述。\par}

{\setlength{\parindent}{17.1bp}（2）对于不符合专利法及其实施细则规定的授予专利权条件的外观设计专利的每项外观设计，均须给出具体的评价说明，并明确结论，必要时应当引证对比文件。\par}

\guidesubsection{gpos.5.10.4.2}{4.2}{专利权评价报告的发送}{}{74.4bp}
专利权评价报告作出后，应当发送给请求人。请求人不是专利权人的，国家知识产权局应当将专利权评价报告出具情况告知专利权人。

\guidesection{gpos.5.10.5}{5.}{专利权评价报告的查阅与复制}{}{74.4bp}
根据专利法实施细则第六十三条的规定，国家知识产权局在作出专利权评价报告后，任何单位或者个人可以查阅或者复制。查阅、复制的相关手续参见本指南第五部分第四章第5.3节的规定。

\guidesection{gpos.5.10.6}{6.}{专利权评价报告的更正}{\guidenote[-11.5bp]{细则 63.2}}{74.4bp}
作出专利权评价报告的部门在发现专利权评价报告中存在错误后，可以自行更正。请求人认为专利权评价报告存在需要更正的错误的，可以请求更正。

更正后的专利权评价报告应当及时发送给请求人。

\guidesubsection{gpos.5.10.6.1}{6.1}{可更正的内容}{}{74.4bp}
专利权评价报告中存在下列错误的，可以进行更正：

{\setlength{\parindent}{17.1bp}（1）著录项目信息或文字错误；\par}

{\setlength{\parindent}{17.1bp}（2）作出专利权评价报告的程序错误；\par}

{\setlength{\parindent}{17.1bp}（3）法律适用明显错误；\par}

{\setlength{\parindent}{17.1bp}（4）结论所依据的事实认定明显错误；\par}

{\setlength{\parindent}{17.1bp}（5）其他应当更正的错误。\par}

\guidesubsection{gpos.5.10.6.2}{6.2}{更正程序的启动}{}{74.4bp}
{\setlength{\parindent}{17.1bp}（1）作出专利权评价报告的部门自行启动\par}

作出专利权评价报告的部门在发现专利权评价报告中存在需要更正的错误后，可以自行启动更正程序。

{\setlength{\parindent}{17.1bp}（2）请求人请求启动\par}

请求人认为作出的专利权评价报告存在需要更正的错误的，可以在收到专利权评价报告后两个月内提出更正请求，请求人不是专利权人的，专利权人可以在上述期限内提出更正请求。期满后提交的，其请求视为未提出。

提出更正请求的，应当以意见陈述书的形式书面提出，写明需要更正的内容及更正的理由，但不得修改专利文件。

\guidesubsection{gpos.5.10.6.3}{6.3}{更正程序的进行和终止}{}{74.4bp}
更正程序启动后，作出专利权评价报告的部门应当成立由组长、主核员和参核员组成的三人复核组，对原专利权评价报告进行复核。复核结果经复核组合议作出，合议时采取少数服从多数的原则。作出原专利权评价报告的审查员和审核员不参加复核组。

复核组认为更正理由不成立，原专利权评价报告无误、不需更正的，应当发出专利权评价报告复核意见通知书，说明不予更正的理由，更正程序终止。

复核组认为更正理由成立，原专利权评价报告有误、确需更正的，应当发出更正的专利权评价报告，并在更正的专利权评价报告上注明以此报告代替原专利权评价报告，更正程序终止。

在更正程序中，复核组一般不进行补充检索，除非因事实认定发生变化，导致原来的检索不完整或者不准确。针对专利权评价报告，一般只允许提出一次更正请求，但对于复核组在补充检索后重新作出的专利权评价报告，请求人可以再次提出更正请求。

\guidechapter{gpos.5.11}{第十一章}{专利开放许可}{}{0bp}
\guidesection{gpos.5.11.1}{1.}{引言}{}{74.4bp}
\guidenote{法 50 及 51}根据专利法第五十条、第五十一条和专利法实施细则第八十五条至第八十八条的规定制定本章。

根据专利法第五十条和第五十一条的规定，专利权人自愿以书面方式向国家知识产权局声明愿意许可任何单位或者个人实施其专利，并明确许可使用费支付方式和标准的，由国家知识产权局予以公告，实行开放许可。任何单位或者个人有意愿实施开放许可专利的，以书面方式通知专利权人，并依照公告的许可使用费支付方式、标准支付许可使用费后，即获得专利实施许可。

提出开放许可声明的专利权人应当声明许可任何单位或者个人在中国境内实施其开放许可的专利。

国家知识产权局对专利权人提出的开放许可声明是否符合规定进行审查后，发出是否准予公告的通知。

本章对专利开放许可声明的提出，专利开放许可声明的撤回，专利开放许可的登记和公告、专利开放许可实施合同的生效，专利开放许可实施合同的备案，专利开放许可实施期间费减手续的办理，以及已实行开放许可专利的相关手续办理做出规定。

\guidesection{gpos.5.11.2}{2.}{开放许可相关原则}{}{74.4bp}
建立专利开放许可制度的目的是为了促进专利技术的实施与运用，通过国家知识产权局公告专利开放许可信息，帮助专利技术供需双方对接。专利开放许可相关程序应当遵循以下原则。

{\setlength{\parindent}{17.1bp}（1）自愿原则\par}

对开放许可声明中的许可条件，在符合相关规定的前提下，当事人可以依照自愿原则设立。

{\setlength{\parindent}{17.1bp}（2）合法原则\par}

为维护开放许可交易安全，国家知识产权局公告开放许可专利的专利权应当有效。已经实行开放许可的专利权有专利法实施细则第八十六条规定情形的，专利权人应当及时撤回开放许可声明，同时通知被许可人。

{\setlength{\parindent}{17.1bp}（3）公开原则\par}

专利开放许可声明予以公告后，国家知识产权局可以通过查阅复制等途径对外公开开放许可声明内容。

\guidesection{gpos.5.11.3}{3.}{专利开放许可声明的提出}{}{74.4bp}
\guidenote{法 50.1}专利权人有意愿实行开放许可专利的，应当向国家知识产权局提交专利开放许可声明。专利开放许可声明原则上应当通过电子形式提交，电子形式提交确有困难的，可以向国家知识产权局指定的地点面交或者邮寄相关文件。

\guidesubsection{gpos.5.11.3.1}{3.1}{专利开放许可声明的客体}{}{74.4bp}
\guidenote{细则 85.1 及 86}专利开放许可声明的客体应当是已经授权公告的发明专利、实用新型专利或者外观设计专利。

专利权有下列情形之一的，专利权人不得对其实行开放许可：

{\setlength{\parindent}{17.1bp}（1）专利权处于独占或者排他许可有效期限内的；\par}

{\setlength{\parindent}{17.1bp}（2）因专利权的归属发生纠纷或者人民法院裁定对专利权采取保全措施，已经中止有关程序的；\par}

{\setlength{\parindent}{17.1bp}（3）没有按照规定缴纳年费的；\par}

{\setlength{\parindent}{17.1bp}（4）专利权被质押，未经质权人同意的；\par}

{\setlength{\parindent}{17.1bp}（5）专利权已经终止的；\par}

{\setlength{\parindent}{17.1bp}（6）专利权已经被宣告全部无效的；\par}

{\setlength{\parindent}{17.1bp}（7）实用新型或者外观设计专利权人未提交专利权评价报告的；\par}

{\setlength{\parindent}{17.1bp}（8）专利权评价报告结论认为实用新型或者外观设计专利权不符合授予专利权条件的；\par}

{\setlength{\parindent}{17.1bp}（9）其他妨碍专利权有效实施的情形。\par}

\guidesubsection{gpos.5.11.3.2}{3.2}{提出专利开放许可声明的主体}{}{74.4bp}
根据专利法实施细则第八十五条第一款的规定，专利权人可以向国家知识产权局提出开放许可声明。共有人就共有专利权提出开放许可声明的，应当提交全体共有人同意的证明材料。针对开放许可相关事宜，有约定的从约定，并提交相关证明材料。

\guidesubsection{gpos.5.11.3.3}{3.3}{专利开放许可声明}{}{74.4bp}
\guidenote{细则 85.2 及}\guidenote[14.5bp]{85.3}专利权人应当按照规定的格式提交专利开放许可声明和其他需要提供的材料。专利权人提交的材料内容应当真实、准确、清楚，符合国家法律规定和社会公德、公共利益的要求，不得出现商业性宣传用语。

专利开放许可声明应当写明以下事项：

{\setlength{\parindent}{17.1bp}（1）专利号；\par}

{\setlength{\parindent}{17.1bp}（2）专利权人的姓名或者名称；\par}

{\setlength{\parindent}{17.1bp}（3）专利许可使用费支付方式、标准；\par}

{\setlength{\parindent}{17.1bp}（4）专利许可期限；\par}

{\setlength{\parindent}{17.1bp}（5）专利权人的联系方式；\par}

{\setlength{\parindent}{17.1bp}（6）专利权人对符合开放许可声明条件的承诺；\par}

{\setlength{\parindent}{17.1bp}（7）其他需要明确的事项。\par}

专利权人应当一并提交对许可使用费计算依据和方式的简要说明，一般不超过2000字。专利许可使用费应当以该简要说明为依据，以固定费用标准支付的，一般不高于2000万元。高于2000万元的，专利权人可以利用专利法第五十条规定的开放许可以外的其它方式进行许可。以提成费支付的，净销售额提成一般不高于20\%，利润额提成一般不高于40\%。

专利开放许可声明应当由专利权人签字或者盖章；专利权属于多个专利权人共有的，可以由代表人签字或者盖章，同时附具共有专利权人签字或者盖章的同意开放许可的书面声明；委托专利代理机构的，专利开放许可声明应当由专利代理机构盖章，同时附具全体专利权人签字或者盖章的同意开放许可的书面声明。

\guidesubsection{gpos.5.11.3.4}{3.4}{准予公告和不予公告}{}{74.4bp}
{\setlength{\parindent}{17.1bp}（1）专利开放许可声明经审查符合规定的，国家知识产权局准予公告。\par}

{\setlength{\parindent}{17.1bp}\guidenote{法 50.1}（2）专利开放许可声明经审查不符合专利法实施细则第八十五条规定、或者属于第八十六条所列情形的，国家知识产权局不予公告。\par}

{\setlength{\parindent}{17.1bp}\guidenote{细则 88}（3）专利权人通过提供虚假材料、隐瞒事实等手段作出开放许可声明的，国家知识产权局一经发现，应予撤销。\par}

\guidesubsection{gpos.5.11.3.5}{3.5}{专利开放许可声明的生效}{}{74.4bp}
专利开放许可声明自公告之日起生效。

\guidesection{gpos.5.11.4}{4.}{专利开放许可声明的撤回}{}{74.4bp}
\guidenote{法 50.2}专利权人可以依据专利法实施细则第八十六条规定或者基于其他正当理由撤回开放许可声明。共有人就共有专利权撤回开放许可声明的，应当取得全体共有人的书面同意。专利权人撤回开放许可声明的，应当提交撤回开放许可声明请求。撤回开放许可声明请求应当由专利权人签字或者盖章；专利权属于多个专利权人共有的，可以由代表人签字或者盖章，同时附具共有专利权人签字或者盖章的同意撤回开放许可的书面声明；委托专利代理机构的，撤回开放许可声明请求应当由专利代理机构盖章，同时附具全体专利权人签字或者盖章的同意撤回开放许可的书面声明。撤回开放许可声明不得附有任何条件。

撤回开放许可声明请求经审查，符合规定的，国家知识产权局准予公告撤回开放许可声明；不符合规定的，国家知识产权局不予公告撤回开放许可声明，并向专利权人说明理由。

开放许可声明的撤回，自公告之日起生效。

\guidenote{细则 86}对于专利权人应当撤回但未及时撤回专利开放许可声明的，国家知识产权局将终止或撤销该专利开放许可声明，并进行公告。

\guidesection{gpos.5.11.5}{5.}{专利开放许可的登记和公告}{}{74.4bp}
\guidenote{法 50}专利开放许可声明有关事项在专利登记簿上登记，并在专利公报上公告。

专利开放许可声明公布的项目包括：主分类号、专利号、开放许可声明编号、专利权人、发明名称、申请日、授权公告日、专利许可使用费支付方式和标准、专利许可期限、专利权人联系方式、开放许可声明生效日等。

撤回专利开放许可声明公布的项目包括：主分类号、专利号、开放许可声明编号、专利权人、发明名称、开放许可声明撤回日等。

\guidesection{gpos.5.11.6}{6.}{专利开放许可实施合同的生效}{}{74.4bp}
\guidenote{法 51.1}任何单位或者个人以书面方式通知专利权人愿意实施其开放许可专利，并依照公告支付许可使用费的，专利开放许可实施合同生效，但有关法律、行政法规另有规定的除外。

中国境内单位或者个人实行专利开放许可的，外国人、外国企业或者外国其他组织有意愿实施时，应当符合《中华人民共和国技术进出口管理条例》和《技术进出口合同登记管理办法》等有关规定。

中国境内单位或者个人实行专利开放许可的，香港、澳门或者台湾地区的个人、企业或者其他组织有意愿实施时，参照上述规定。

\guidesection{gpos.5.11.7}{7.}{专利开放许可实施合同的备案}{}{74.4bp}
\guidenote{细则 87}专利权人或者被许可人应当在开放许可实施合同生效后，凭能够证明达成开放许可的书面文件向国家知识产权局办理备案手续。

办理专利开放许可实施合同备案的应当提交下列文件：

{\setlength{\parindent}{17.1bp}（1）请求人签章的专利实施许可合同备案申请表；\par}

{\setlength{\parindent}{17.1bp}（2）被许可人以书面方式向专利权人发出的通知；\par}

{\setlength{\parindent}{17.1bp}（3）被许可人向专利权人支付许可使用费的凭证（或专利权人收到许可使用费的凭证）；\par}

{\setlength{\parindent}{17.1bp}（4）请求人身份证明；\par}

{\setlength{\parindent}{17.1bp}（5）委托代理的，注明委托权限的委托书；\par}

{\setlength{\parindent}{17.1bp}（6）经办人身份证明；\par}

{\setlength{\parindent}{17.1bp}（7）其他需要提供的材料。\par}

专利开放许可实施合同备案手续的办理参照《专利实施许可合同备案办法》执行。

\guidesection{gpos.5.11.8}{8.}{专利开放许可实施期间费减手续的办理}{}{74.4bp}
\guidenote{法 51.2}专利开放许可实施期间是指专利开放许可实施合同生效之日至专利许可期限届满期间。

请求人办理专利开放许可实施合同备案，视为专利权人同时提出专利年费减缴请求。专利开放许可实施合同被准予备案的，专利权人可以在专利开放许可实施期间，按规定享有自备案日起尚未到期的专利年费的减缴。专利权人撤回开放许可声明的，自下一专利年度起不再享有因开放许可获得的专利年费减缴。专利权人同时符合两项专利年费减缴条件的，按照其中减缴比例较高的一项条件予以减缴。

\guidenote{法 51.3}实行开放许可的专利权人与被许可人就许可使用费进行协商后签订普通许可合同的，不属于开放许可。

\guidesection{gpos.5.11.9}{9.}{已实行开放许可的专利相关手续办理}{}{74.4bp}
对于已实行开放许可的专利，在办理以下手续前，专利权人应当首先撤回开放许可声明：

{\setlength{\parindent}{17.1bp}（1）因专利权转让，提出著录项目变更请求的；\par}

{\setlength{\parindent}{17.1bp}（2）专利权人以书面声明放弃其专利权的。\par}

除专利权转让外，专利权人因其他事由发生变更且继续实行开放许可的，应当及时办理原开放许可声明撤回和重新声明的相关手续；专利权人变更后不再实行开放许可的，应当及时办理撤回原开放许可声明的手续。

专利权人以实行开放许可的专利权出质的，办理专利权质押登记时，应当提供质权人同意继续实行专利开放许可的书面声明。

\stopparttoc
